\documentclass[11pt]{amsart}
\usepackage{amsmath,amssymb}
\ifdefined\XeTeXrevision
  \IfFileExists{lmroman7-regular.otf}{}{%
    \DeclareFontFamily{TU}{paperlmr}{}%
    \DeclareFontShape{TU}{paperlmr}{m}{n}{<->\UnicodeFontFile{lmroman10-regular}{\UnicodeFontTeXLigatures}}{}%
    \DeclareFontShape{TU}{paperlmr}{m}{it}{<->\UnicodeFontFile{lmroman10-italic}{\UnicodeFontTeXLigatures}}{}%
    \DeclareFontShape{TU}{paperlmr}{m}{sl}{<->ssub*paperlmr/m/it}{}%
    \DeclareFontShape{TU}{paperlmr}{m}{sc}{<->\UnicodeFontFile{lmromancaps10-regular}{\UnicodeFontTeXLigatures}}{}%
    \DeclareFontShape{TU}{paperlmr}{b}{n}{<->\UnicodeFontFile{lmroman10-bold}{\UnicodeFontTeXLigatures}}{}%
    \DeclareFontShape{TU}{paperlmr}{bx}{n}{<->ssub*paperlmr/b/n}{}%
    \DeclareFontFamily{TU}{paperlmtt}{}%
    \DeclareFontShape{TU}{paperlmtt}{m}{n}{<->\UnicodeFontFile{lmmono10-regular}{}}{}%
    \renewcommand{\rmdefault}{paperlmr}%
    \renewcommand{\ttdefault}{paperlmtt}%
    \DeclareMathAlphabet{\mathsf}{OT1}{cmr}{m}{n}%
    \DeclareFontFamily{U}{msa}{}%
    \DeclareFontFamily{U}{msb}{}%
    \DeclareFontShape{U}{msa}{m}{n}{<->msam10}{}%
    \DeclareFontShape{U}{msb}{m}{n}{<->msbm10}{}%
  }
\fi
\usepackage[hidelinks,pdfusetitle]{hyperref}
\numberwithin{equation}{section}
\newtheorem{theorem}{Theorem}[section]
\newtheorem{proposition}[theorem]{Proposition}
\newtheorem{lemma}[theorem]{Lemma}
\newtheorem{corollary}[theorem]{Corollary}
\theoremstyle{remark}
\newtheorem{remark}[theorem]{Remark}
\newcommand{\Z}{\mathbb Z}
\newcommand{\R}{\mathbb R}
\newcommand{\C}{\mathbb C}
\newcommand{\F}{\mathbb F}
\newcommand{\T}{\mathbb T}
\newcommand{\E}{\mathbb E}
\newcommand{\one}{\mathbf 1}
\newcommand{\Typ}{\mathsf{Typ}}
\newcommand{\supp}{\operatorname{supp}}
\newcommand{\Var}{\operatorname{Var}}
\newcommand{\SL}{\operatorname{SL}}
\newcommand{\BV}{\mathrm{BV}}
\setlength{\emergencystretch}{2em}

\title[Reflected binary Chowla correlations under Dirichlet GRH]{Reflected Binary Chowla Correlations for All Integers under Dirichlet GRH}
\hypersetup{pdftitle={Reflected Binary Chowla Correlations for All Integers under Dirichlet GRH}}
\author{KenRui Lin}
\address{Independent Researcher, Guangzhou, China}
\email{longinuslkr@gmail.com}
\subjclass[2020]{11N37, 11N25, 11L05, 11M26}
\keywords{M\"obius function, reflected Chowla correlation, Dirichlet GRH, Kloosterman sums, Type I and Type II sums}

\begin{document}
\begin{abstract}
Assume the generalized Riemann hypothesis for Dirichlet $L$-functions. For every fixed complex-valued weight of bounded variation on $[0,1]$, we prove that the weighted correlation of $\mu(a)$ and $\mu(N-a)$ is $o(N)$ as $N$ ranges over all positive integers. Exact greatest-common-divisor slices give two unit coordinates at the physical endpoints of a determinant pairing. A Kloosterman count with both unit constraints averages signed path moments on the resulting index-$Q$ lattices, where $Q$ is the reduced modulus. On the minor arcs we bound these moments with a history operator and distinguish collisions by partition rank. On the major arcs a modulus lift separates the induced characters for which the label cancellation can fail. The underlying reduced pair estimate gives an arbitrary fixed logarithmic saving with a suitable typicality sieve.
\end{abstract}
\maketitle

\section{Introduction and statements}

Write $\mu$ for the M\"obius function, $L=\log N$ and $e(t)=e^{2\pi i t}$. Dirichlet GRH means the generalized Riemann hypothesis for primitive Dirichlet $L$-functions, including the Riemann zeta function. The reflected correlation is
\[
 C_N(w)=\sum_{1\le a<N}\mu(a)\mu(N-a)w(a/N).
\]
The parameter $N$ serves both as the reflection point and, after taking greatest-common-divisor slices, as the source of the determinant modulus. We build on the decomposition and determinant method of Lichtman \cite{Lichtman}, with particular attention to the unit conditions and shared prime powers that occur for general moduli.

\begin{theorem}[Reflected binary Chowla for all integers]\label{thm:main}
Assume Dirichlet GRH. For every fixed $w\in \BV([0,1];\C)$,
\[
 C_N(w)=o_w(N)\qquad(N\longrightarrow\infty,\ N\in\Z_{\ge2}).
\]
The limit is uniform over the integer parameter, and hence holds on the even integers, the primorials, the powers of two, the semiprimes, and the primes.
\end{theorem}

For the pair estimate put $\alpha=1/1000$ and $\beta=1/100$. Set
\[
 \begin{gathered}
 b_{\rm pat}=\left\lfloor\frac{(\beta-\alpha)\log L}{\log 2}\right\rfloor,
 \qquad \lambda_j^0=2^{j-1}L^\alpha,\\
 \mathcal P_j^0=\{\ell\text{ prime}:e^{\lambda_j^0}<\ell\le e^{2\lambda_j^0},\ \ell\nmid N\}
 \quad(1\le j\le b_{\rm pat}).
 \end{gathered}
\]
For a fixed positive integer $k$, $\Typ(v)$ is the indicator that $v$ has a prime divisor in at least $3k$ of the sets $\mathcal P_j^0$.

\begin{theorem}[Uniform log-power pair bound on reduced moduli]\label{thm:pair}
Assume Dirichlet GRH and fix $K>2$, $C_g\ge1$ and $A>0$. There is a fixed $k=k(A,K,C_g)$ with the following property. Let $g\mid N$ be squarefree, $g\le L^{C_g}$ and $Q=N/g$. Let $L^2\le r_1<r_2\le L^K$ be primes with $r_1,r_2\nmid N$, and set
\[
 U=Q-r_1j,\qquad V=Q-r_2j,\qquad
 W_w(j)=w(r_1j/Q)\overline{w(r_2j/Q)}.
\]
Uniformly over these choices and real intervals $\mathcal I$,
\begin{equation}\label{eq:pair}
 \left|\sum_{\substack{j\in\mathcal I\cap\Z\\1\le j<Q/r_2}}
 \mu(U)\mu(V)\one_{(UV,N)=1}\Typ(U)\Typ(V)W_w(j)\right|
 \ll_{A,K,C_g,w} QL^{-A}.
\end{equation}
Taking $A=2K+4$ gives $O_{K,C_g,w}(NL^{-2K-4})$ on each slice.
\end{theorem}

\begin{remark}[Bounded variation and the Goldbach remainder]\label{rem:scope}
The bounded-variation hypothesis cannot be replaced by bounded Borel measurability. Define a single bounded Borel function by setting $w(a/p)=\mu(a)\mu(p-a)$ for each prime $p$ and $1\le a<p$, and setting $w=0$ elsewhere. The reduced denominator makes this assignment unambiguous. Then $C_p(w)$ counts the $a$ for which both $a$ and $p-a$ are squarefree. A union bound for the two square divisibility conditions gives
\[
 C_p(w)\ge p-1-2p\sum_{\ell\text{ prime}}\ell^{-2}\gg p.
\]
Indeed, its sum over primes is at most $1/4+\sum_{n\ge3,\ n\ \mathrm{odd}}n^{-2}=\pi^2/8-3/4<1/2$.

For the even Goldbach problem, the identity $\Lambda(n)=-\sum_{d\mid n}\mu(d)\log d$ leads to the long-divisor remainder $R_{Z_{\rm G}}(n)=-\sum_{d\mid n,\ d>Z_{\rm G}}\mu(d)\log d$. Its additive self-convolution can be written as
\[
 (R_{Z_{\rm G}}*_{\rm add} R_{Z_{\rm G}})(N)=
 \sum_{u,v\ge1}\ \sum_{\substack{ua+vb=N\\a,b>Z_{\rm G}}}
 \mu(a)\mu(b)\log a\log b.
\]
Theorem~\ref{thm:main} gives cancellation for the underlying reflected kernel at $u=v=1$. Theorem~\ref{thm:pair} gives cancellation for the typical reduced pairs at logarithmic-power dilation scales. The intermediate power scales from $N^\varepsilon$ up to the $N^{7/15}$ range, and the summation with products of logarithmic weights at the scale of the Goldbach main term, remain to be controlled in the full remainder. Here $7/15$ refers to a balanced right-end threshold in the separate trilinear approach using the spectral Kloosterman estimates of \cite[Corollary~5.5]{Pascadi}.
\end{remark}

\paragraph{Analytic inputs.}
We use the quantitative prime number theorem \cite[Theorem~1.1 and Corollary~1.2]{JohnstonYang}, the Dirichlet GRH explicit formula \cite[Lemma~2]{KandhilLanguascoMoree}, and the full-modulus Weil bound in \cite[Eq.~(1.2)]{BlomerMan}. The numbered references to Lichtman are to arXiv:2611.28392v1. Its Proposition~5.2 and Eq.~(5.8) supply the Type I/II decomposition, with Lemmas~5.1 and~5.3 controlling its truncations. Lemma~5.4 supplies a divisor-weighted M\"obius Bombieri--Vinogradov estimate for all residue classes. We give the reflected versions of the Type I sieve of Proposition~5.5 and of the determinant preparation of Lemmas~6.1--6.2. Proposition~8.1 supplies the rough M\"obius character estimate. The geometric and history argument follows the framework of \cite[Sections~9--11]{Lichtman} and the lattice-box reduction of \cite[Section~4.4]{LuWangWu}. We give the details of these steps for the present correlation.

\section{Typicality, exact gcd slices, and the K\'atai reduction}\label{sec:katai}

We use $\tau_j$ for the $j$-fold divisor function and $\tau=\tau_2$. The quantitative prime number theorem and partial summation give, uniformly for $L^\alpha\le\lambda\le L^\beta$,
\[
 \sum_{e^\lambda<\ell\le e^{2\lambda}}\ell^{-1}=\log 2+O_T(L^{-T})
 \quad (T>0).
\]
The loss on excluding primes dividing $N$ is at most $\omega(N)e^{-\lambda}\ll Le^{-\lambda}$, where $\omega(N)$ counts distinct prime factors. Since $\sum_{\ell>e^\lambda}\ell^{-2}\ll e^{-\lambda}$, each interval also satisfies
\begin{equation}\label{eq:prime-density}
 \prod_{\ell\in\mathcal P_j^0}(1-\ell^{-1})=\tfrac12+O_T(L^{-T}).
\end{equation}

For clarity, a sieve truncation used below is the identity for the no-hit indicator up to an error bounded by
\begin{equation}\label{eq:bonf}
 \left|\one_{(v,\prod_{\ell\in\mathcal H}\ell)=1}
 -\sum_{\substack{d\mid v,\ d\mid\prod_{\ell\in\mathcal H}\ell\\\omega(d)\le h}}\mu(d)\right|
 \le \sum_{\substack{d\mid v,\ d\mid\prod_{\ell\in\mathcal H}\ell\\\omega(d)=h+1}}1 .
\end{equation}
This follows from the alternating binomial sum at $\omega((v,\prod_{\ell\in\mathcal H}\ell))$. If $\sum_{\ell\in\mathcal H}1/\ell\ll\log L$, the reciprocal mass of divisors with $h+1$ prime factors is at most $(\sum_{\ell\in\mathcal H}1/\ell)^{h+1}/(h+1)!$. With $h=\lceil c\log L\rceil$, this is $O_T(L^{-T})$ for sufficiently large fixed $c=c(T)$, by the factorial bound. The corresponding tail of the expansion of the density product has the same bound up to an absolute factor. When the primes are at most $\exp(O(L^\zeta))$ for fixed $\zeta<1$, all divisors in the truncated sums are $N^{o(1)}$, so an $O(1)$ counting error for each divisor costs $N^{o(1)}$.

For a non-typical $v\le N$, at most $3k-1$ intervals are hit. Fix a set of $3k-1$ exempt intervals and sieve against all the others, counting multiples of each sieve divisor by $N/d+O(1)$. Equation~\eqref{eq:prime-density} and the preceding truncation give
\begin{equation}\label{eq:atyp}
 \#\{1\le v\le N:\Typ(v)=0\}
 \ll_k N 2^{-b_{\rm pat}}b_{\rm pat}^{3k-1}+o_k(N)=o_k(N),
\end{equation}
uniformly over $N$ for each fixed $k$.

Let $a$ be uniform modulo $N$, with $\gcd(0,N)=N$. For $d\mid N$,
\begin{equation}\label{eq:gcd-markov}
 \begin{gathered}
 \Pr((a,N)=d)=\frac{\varphi(N/d)}{N},\qquad
 \E\log(a,N)=\sum_{\ell^j\mid N}\frac{\log\ell}{\ell^j}\ll\log L,\\
 \Pr((a,N)>L^{C_g})\ll C_g^{-1}.
 \end{gathered}
\end{equation}
The first identity counts reduced residues. The second uses $\Pr(\ell^j\mid a)=\ell^{-j}$. For $\ell\le L^2$, the geometric sums and partial summation give $\sum_{\ell\le L^2}(\log\ell)/\ell\ll\log L$. For $\ell>L^2$, their contribution is $\ll L^{-2}\sum_{\ell\mid N}\log\ell\ll L^{-1}$. Markov's inequality then gives the tail bound. For primorials $N=\prod_{\ell\le y}\ell$ this tail need not tend to zero at fixed $C_g$. In $(y^{1/2},y]$ choose $h_1>2C_g$ disjoint bands $(y^{c_i},y^{d_i}]$ with fixed endpoints in the exponent. The events of hitting a prime factor in each band are independent by CRT and each has limiting probability $1-c_i/d_i>0$. On their intersection $(a,N)>y^{h_1/2}>L^{C_g}$ for large $y$, since $L\asymp y$. The divisor measure assigning mass proportional to $1/d$ to each $d\mid N$ also has a positive limiting lower bound on the tail by the same argument, with independent inclusion probability $1/(\ell+1)$ and total mass $\prod_{\ell\le y}(1+1/\ell)\asymp\log y$ before normalization.

Set $R=L^K$, $\mathcal P=\{r\text{ prime}:L^2\le r\le R,\ r\nmid N\}$, and $\sigma=\sum_{r\in\mathcal P}1/r$. Partial summation and the prime number theorem give $\sigma=\log(K/2)+o(1)$, because the omitted mass is $\ll\omega(N)L^{-2}$. Set
\[
 F(a)=\mu(N-a)w(a/N)\one_{(a,N)\le L^{C_g}}\Typ(N-a)
 \quad(1\le a<N),
\]
with $F=0$ off this range. By \eqref{eq:atyp}--\eqref{eq:gcd-markov}, replacing $C_N(w)$ by $\sum_a\mu(a)F(a)$ costs $O_w(N/C_g)+o_{k,w}(N)$. These are errors on the original sum, before pair formation.

Expanding the square and counting multiples of $r$ and $r r'$ gives
\[
 \sum_{1\le a<N}\left|\sum_{r\in\mathcal P}\one_{r\mid a}-\sigma\right|^2
 \le N\sigma+O(R^2).
\]
Cauchy--Schwarz compares $\sigma\sum_a\mu(a)F(a)$ with $\sum_{r,s}\mu(rs)F(rs)$ at cost $\ll_w\sqrt{N(N\sigma+R^2)}$. Replacing $\mu(rs)$ by $-\mu(s)$ costs $\ll_w N\sum_{r\in\mathcal P}r^{-2}=o_w(N)$, since the error is supported on $r\mid s$. One more Cauchy--Schwarz application bounds the square of the resulting sum by
\begin{equation}\label{eq:katai}
 N\left(O_w(N\sigma)+2\sum_{r_1<r_2\in\mathcal P}
       \left|\sum_s F(r_1s)\overline{F(r_2s)}\right|\right).
\end{equation}

For a pair in \eqref{eq:katai}, let $g=(s,N)$. As $r_i\nmid N$, all of $(r_i s,N)$ and $(N-r_i s,N)$ equal $g$. Neither M\"obius factor can survive a non-squarefree $g$. For squarefree $g$ and $s=gj$, one has $N-r_1s=gU$, $N-r_2s=gV$, and
\begin{equation}\label{eq:gcd-slice}
 \mu(gU)\mu(gV)=\mu(U)\mu(V)\one_{(UV,g)=1},\qquad
 (j,Q)=1\quad\Longleftrightarrow\quad(UV,Q)=1.
\end{equation}
Every sieve prime is coprime to $g$, so $\Typ(gU)=\Typ(U)$ and $\Typ(gV)=\Typ(V)$. Also $(UV,Q)=(UV,g)=1$ is equivalent to $(UV,N)=1$, including when $(g,Q)>1$. Each cross term is therefore a sum of exactly the slices in Theorem~\ref{thm:pair}. Summing the bound $QL^{-A}$ over $g\le L^{C_g}$ costs at most $NL^{-A+C_g}$. Take $A=2K+4+C_g$ and use its associated fixed $k$. All cross terms together contribute $O_{K,C_g,w}(NL^{-4})$ in the parentheses in \eqref{eq:katai}. Dividing by $N\sigma$ after taking the square root and recovering the two truncations yields
\begin{equation}\label{eq:limsup}
 \limsup_{N\to\infty}\frac{|C_N(w)|}{N}
 \ll_w \frac{1}{C_g}+\frac{1}{\sqrt{\log(K/2)}}.
\end{equation}
The implied constant here is independent of $K,C_g$. Letting these fixed parameters grow proves Theorem~\ref{thm:main} once Theorem~\ref{thm:pair} is established.

\section{From the reduced pair to Type I and Type II}\label{sec:types}

Fix $g,r_1,r_2$ as in Theorem~\ref{thm:pair}. Write $\Delta=r_2-r_1$. On $1\le j<Q/r_2$, the two reduced arguments satisfy
\begin{equation}\label{eq:reflection}
 V>0,\qquad Q/R\le U<Q,\qquad r_2U-\Delta Q=r_1V.
\end{equation}
It suffices to bound the unweighted sum on each subinterval. Indeed,
$\|W_w\|_\infty\le\|w\|_\infty^2$ and $\Var(W_w)\le2\|w\|_\infty\Var(w)$, because the two arguments are monotone maps into $[0,1]$. Discrete partial summation then gives the desired weighted bound with a constant depending only on $w$ in this step.

For $V$ with $\Typ(V)=1$, record its earliest $3k$ hit intervals. This gives $O_k(b_{\rm pat}^{3k})=L^{o(1)}$ patterns. The pattern indicator $\phi$ requires hits in those intervals and misses in the intervening unselected intervals up to the last hit. Make the first, fourth, and every third hit thereafter active, leaving all remaining conditions in $\phi^-$. Thus $\phi=\phi^+\phi^-$, where $\phi^+$ tests the $k$ active hits. Multiplication by an active prime leaves $\phi^-$ unchanged. Renumber the active prime sets as $\mathcal P_i$, with logarithmic endpoints $\lambda_i$. Then $\lambda_1\ge L^\alpha$, $\lambda_k\ll L^\beta$ and $\lambda_{i+1}\ge8\lambda_i$. The primes in the patterns are all coprime to $\Delta N r_1r_2$, since they avoid $N$ and exceed $R$ for every sufficiently large $N$. This remains true of all products of active primes used later.

Put $W=N^{1/100}$, $z=\exp(L^\zeta)$ and $\zeta=(1+\beta)/2$. Denote the indicator of integers all of whose prime factors exceed $z$ by $R_z$. Write $u=s_u t_u$ with $s_u$ supported on primes $\le z$ and $t_u$ on primes $>z$. To apply \cite[Proposition~5.2]{Lichtman} at the original parameter $N$, we use its convolution identity throughout the actual range $U\le Q$. Let
\[
 \begin{gathered}
 \nu(u)=\mu(u)R_z(u),\qquad
 a_0=\nu\one_{u\le W^2},\qquad b_0=\nu\one_{u>W^2},\\
 f_z(u)=\mu(u)\one_{P^+(u)\le z}\one_{u\le W},
 \end{gathered}
\]
with $P^+(1)=1$. The function $F_z=\nu*f_z$ equals $\mu(u)\one_{s_u\le W}$ and satisfies
\begin{equation}\label{eq:convolution}
 F_z=2a_0*f_z - A*R_z+b_0*\mathcal B,
 \qquad A=a_0*a_0*f_z,\qquad \mathcal B=b_0*R_z*f_z.
\end{equation}
Here $*$ is Dirichlet convolution. In fact $\nu*R_z=\one_{u=1}$, whence $\nu=(a_0+b_0)*(a_0+b_0)*R_z$ and $a_0=a_0*(a_0+b_0)*R_z$. Substitution proves \eqref{eq:convolution}. The decomposition into smooth and rough factors gives $|A(q_{\mathrm I})|\le\tau(q_{\mathrm I})$, supported on $q_{\mathrm I}\le W^5$, and $|\mathcal B(t)|\le\tau(t)$, supported on $t>W^2$.

The cost of replacing $\mu(U)$ by $F_z(U)$ is $O_D(QL^{-D})$ for every fixed $D$. To see this, count $U\le Q$ with $s_U>W$ by $Q\sum_{s>W,\ P^+(s)\le z}1/s$. With $\epsilon_z=1/\log z$, the latter reciprocal sum is bounded by
\[
 W^{-\epsilon_z}\prod_{\ell\le z}(1-\ell^{-1+\epsilon_z})^{-1}
 \ll L^{O(1)}\exp(-L^{1-\zeta}/100).
\]
The product bound follows by taking logarithms and using $\ell^{\epsilon_z}\le e$ and the reciprocal-prime estimate. The map $j\mapsto U$ is injective, so this estimate holds with the original indicator restrictions. This is the approximation in \cite[Lemma~5.1]{Lichtman} in our range. The support of $a_0*f_z$ is $U\le W^3$, disjoint from $U\ge Q/R$ because $Q\ge N L^{-C_g}$. The last two terms of \eqref{eq:convolution} give Type I and Type II. For Type II, truncating $\mathcal B(t)$ to $|\mathcal B(t)|\le L^C$ changes the sum by at most
\begin{equation}\label{eq:coefficient-trunc}
 QL^{-C}\sum_{t\le Q}\frac{\tau(t)^2}{t}\ll QL^{4-C},
\end{equation}
as in \cite[Lemma~5.3]{Lichtman}. We used $\tau^2\le\tau_4$ and the elementary divisor-sum bound $\sum_{t\le x}\tau_j(t)/t\ll_j(\log(2x))^j$, obtained by summing over ordered factors. The related bound $\sum_{t\le x}\tau_j(t)\ll_j x(\log(2x))^{j-1}$ follows by counting the last factor.

\subsection{Type I on all residue classes}\label{sec:type-I}

Write $U=q_{\mathrm I}h$, with $q_{\mathrm I}\le W^5$ and $R_z(h)=1$. All pattern primes are at most $z$, so $\Typ(U)=\Typ(q_{\mathrm I})$ on this support. We restrict $(q_{\mathrm I},N)=1$. By \eqref{eq:gcd-slice}, $(U,Q)=1$ and $(V,Q)=1$ are equivalent. The mask $(h,N)=1$ thus accounts for $(U,Q)=1$ as well as $(U,g)=1$. Temporarily dropping this mask costs at most
\begin{equation}\label{eq:typeI-unit}
 \sum_{\substack{\ell\mid N\\\ell>z}}\sum_{q_{\mathrm I}\le W^5}
 \tau(q_{\mathrm I})\left(\frac{Q}{q_{\mathrm I}\ell}+1\right)
 \ll QL^{O(1)}z^{-1}+Q^{1/20+o(1)}.
\end{equation}
The condition $(V,g)=1$ stays in the sum. We have used $\omega(N)\ll L$ and $W^5=Q^{1/20+o(1)}$.

Use \eqref{eq:bonf} to truncate the roughness test to squarefree divisors $d_0\mid h$ with at most $h_0=\lceil c\log L\rceil$ prime factors. For the pattern on $V$, expand the $3k$ required hits as products of $1$ minus a no-hit test. This is a sum of at most $2^{3k}$ signed no-hit tests, one for each resulting union of prime sets, including the required misses. Apply \eqref{eq:bonf} to each such union, obtaining divisors $d_v\mid V$ with at most $h_0$ prime factors. The two truncations have reciprocal remainders $O_D(L^{-D})$ for any prescribed $D$, on choosing $c$ large. More explicitly, for the first truncation one can retain the actual pattern bounded by one. For the second, the total reciprocal mass of the retained roughness divisors is at most $\prod_{\ell\le z}(1+1/\ell)=L^{O(1)}$. In counting a divisor pair, solvability implies $(d_v,r_2q_{\mathrm I}d_0)=1$, since a common prime would divide $\Delta Q$ by \eqref{eq:reflection}. Therefore, with $h=d_0 t$, the condition $d_v\mid V$ places $t$ in at most one residue class modulo $d_v$, even in the presence of the original $r_1$-condition. Such counts are $O(Q/(q_{\mathrm I}d_0d_v)+1)$. The first term costs $\ll L^2$ from the $q_{\mathrm I}$ reciprocal weights, in addition to the stated sieve masses. The aggregate $+1$ cost is $Q^{1/20+o(1)}$. This proves the required reflected version of the truncation in \cite[Proposition~5.5]{Lichtman}.

For the remaining progression, the physical $U\equiv Q\pmod{r_1}$ excludes $r_1\mid q_{\mathrm I}d_0$. The identity
\[
 r_1V+\Delta Q=r_2q_{\mathrm I}d_0t
\]
therefore fixes $V$ in one class modulo $r_2q_{\mathrm I}d_0$. Intersect with $d_v\mid V$ and expand $\one_{(V,g)=1}=\sum_{e\mid(V,g)}\mu(e)$. Each progression is a single class, or is empty, modulo
\[
 m_0=\operatorname{lcm}(r_2q_{\mathrm I}d_0d_v,e)\le Q^{1/20+o(1)}.
\]
For fixed $e,m_0$ there are at most $\tau(e)$ possibilities for $r_2q_{\mathrm I}d_0d_v$, since each has the form $m_0/d$ with $d\mid e$. At a fixed original modulus the coefficient sum in absolute value is $O_k(\tau_4(m_0))$, by grouping triples with weight $\tau(q_{\mathrm I})$ and using monotonicity of $\tau_4$ under divisibility. Thus the total coefficient is $O_k(\tau(g)^2\tau_4(m_0))$. We use the estimate of \cite[Lemma~5.4]{Lichtman}, in the form
\[
 \begin{gathered}
 \sum_{m_0\le H_*}\tau_4(m_0)
 \max_{b\bmod m_0}\sup_{x\le Q}
 \left|\sum_{\substack{n\le x\\n\equiv b\ (m_0)}}\mu(n)\right|
 \ll_D QL^{-D},\\
 H_*\le Q^{1/2}(\log Q)^{-c_D}.
 \end{gathered}
\]
An interval contributes at most twice the supremum, and $\tau(g)^2\le L^{2C_g}$. By increasing the precision in this estimate and in the sieve remainders, the Type I contribution per pattern is $O_D(QL^{-D})$ for every fixed $D$.

\subsection{Type II boxes and label extraction}\label{sec:type-II}

Absorb $r_2$ into the second variable and write $r_2U=mn$. The first coefficient is a rough M\"obius coefficient with $(m,N)=1$. The second includes $(n,N)=1$, $r_2\mid n$, the truncated $\mathcal B(n/r_2)$, and $\Typ(n/r_2)$. This last test equals $\Typ(U)$ since $m$ is rough. The relation $mn\equiv r_2Q\pmod{r_1}$ lets us split into at most $r_1-1\le R$ blocks indexed by $m\equiv i\pmod{r_1}$, with $n\equiv r_2Qi^{-1}\pmod{r_1}$. Both factors in the decomposition before absorbing $r_2$ exceed $W^2$. In a dyadic box $(M,2M]\times(B_0,2B_0]$ that meets the physical range, we may take $M,B_0\ge N^\delta$ with $\delta=1/200$. Its product length $X=MB_0$ satisfies $\Delta Q/4\le X\ll RQ$. The $O(L^2)$ boxes meet $V>0$. The condition $(V,Q)=1$ is automatic on the $(mn,Q)=1$ support, by \eqref{eq:reflection} and the unit condition on $r_1,r_2$. We keep $(V,g)=1$.

The original condition $j\in\mathcal I$, $1\le j<Q/r_2$ gives an interval for the integer $mn$. Use half-integer endpoints $b_-,b_+$ and the truncated Perron kernel for $(b_\pm/(mn))^s$ on $\operatorname{Re}s=1/L$ at height $N^4$. The error at each endpoint is bounded by a constant times
\[
 \left(\frac{b_\pm}{mn}\right)^{1/L}
 \min\left(1,\frac{1}{N^4|\log(b_\pm/(mn))|}\right).
\]
In nonempty boxes the endpoints may be taken within $[X-1,4X+1]$ and each is a distance at least $1/2$ from any integer product. Consequently the total error is a power saving, even with both coefficients bounded by $O(L^C)$. The separated weights have normalized real powers of bounded total variation and polynomially bounded imaginary exponents. The Perron integrals have $L^1$ cost $O(L)$, from integrating $1/|1/L+it|$. After this separation the positive cofactor condition continues to enforce $V>0$. The first coefficient now has the form
\begin{equation}\label{eq:first-coefficient}
 \begin{gathered}
 \alpha_m=\mu(m)R_z(m)\one_{(m,N)=1}\,m^{it_0}w_m(m),\\
 |t_0|\le N^{O(1)},\qquad
 \|w_m\|_\infty+\Var(w_m)\ll1,
 \end{gathered}
\end{equation}
with the block condition separate. The harmless constants in the normalization of the powers are bounded uniformly over the box and endpoint.

If an active prime $\ell$ has $\ell^2\mid V$, then for fixed rough $m$ the relation $r_1V=mn-\Delta Q$ confines $n$ to a single class modulo $\ell^2$, since $\ell\le z$ and $(m,\ell)=1$. Each box has $O(X/\ell^2)$ such pairs because $\ell^2=N^{o(1)}$ and $B_0\ge N^\delta$. Summing over $\ell\ge e^{\lambda_1}$ gives an error $O(X L^{2C}e^{-\lambda_1})$ by absolute values. On the complementary set, expand the $k$ active hits using the Ramar\'e weights of \cite[Lemma~6.1]{Lichtman}. Let $\mathcal A$ contain products of one prime from each active interval, and set
\[
 c(q_0)=\frac{\mu(q_0)\phi^-(q_0)\one_{(q_0,g)=1}}
               {\prod_{i=1}^k(\omega_{\mathcal P_i}(q_0)+1)},
 \qquad \omega_{\mathcal P_i}(q_0)=\#\{\ell\in\mathcal P_i:\ell\mid q_0\} .
\]
Then $\mu(V)\phi(V)\one_{(V,g)=1}=(-1)^k\sum_{a\mid V,\ a\in\mathcal A}c(V/a)$ in the absence of active squares. Indeed all chosen primes divide $V$ exactly once, $\mu(V)=(-1)^k\mu(V/a)$, and the denominator is the number of choices of $a$. The mask and passive factor are invariant under insertion of $a$. Even on the discarded squareful set the absolute value of the expansion is at most one, because $\omega_{\mathcal P_i}(V/a)+1\ge\omega_{\mathcal P_i}(V)$ for an admissible choice. Thus allowing these terms again costs the same error.

Choose a fixed nonincreasing smooth cutoff $\vartheta$ on $(0,\infty)$, equal to one up to $1$ and zero from $2$ onward. Set $\eta(t)=\vartheta(t/2)-\vartheta(t)$ and $\eta_Y(t)=\eta(t/Y)$ for dyadic $Y$. Then $\eta$ is a nonnegative smooth function supported on $[1,4]$, with $\sum_{Y\in 2^\Z}\eta_Y(t)=1$ for $t>0$. Its normalized dilates have a uniform bound for total variation. There are $O(L)$ relevant dyadic scales because all label sizes are $\exp(O_k(L^\beta))$. At one scale the boxed sum is a sum over $q_0>0$ with weight $(-1)^kc(q_0)$ and factorization $mn-\Delta Q=ar_1q_0$, with $a$ weighted by $\eta_Y(a)$.

For comparison with the unreduced coordinates, if $n_{\rm old}=gn$ then $mn_{\rm old}-\Delta N=g(mn-\Delta Q)$, and the same identity holds for $s_{\rm old}=gs$. As $(ar_1,g)=(br_1,g)=1$, an original common cofactor is necessarily $gq_0$. The reduced equations use $q_0$ itself. None of these identities requires $(\Delta,Q)=1$.

\section{The physical pairing and its deletions}\label{sec:pairing}

Fix one Type II box $(M,2M]\times(B_0,2B_0]$, with
$X=MB_0$, both sides at least $N^\delta$, and
$\Delta Q/4\le X\ll RQ$. Its two coefficients are $\alpha_m$ and
$\beta_n$, of size $O(L^C)$, with $\alpha_m$ supported on $z$-rough
integers. The two coefficients include the masks $(m,N)=(n,N)=1$.
We work in one block $i\in(\mathbb Z/r_1\mathbb Z)^\times$ with
\begin{equation}\label{eq:physblock}
 m,r\equiv i,\qquad n,s\equiv r_2 Q i^{-1}\pmod{r_1},
 \qquad n,s\equiv0\pmod{r_2}.
\end{equation}

Here $r,s$ denote the second pair of box coordinates. The active prime
intervals $\mathcal P_j=(\exp(\lambda_j),\exp(2\lambda_j)]\cap\mathbb P$
are understood to omit primes dividing $N$. They satisfy
$\lambda_1\gg L^\alpha$, $\lambda_k\ll_k L^\beta$, and
$\lambda_{j+1}\ge8\lambda_j$. Write $\mathcal A$ for the products of
one prime from each interval, $V_j=\sum_{\ell\in\mathcal P_j}1/\ell$,
and $\mathcal V=\prod_j V_j$, so $V_j\asymp 1$. For a dyadic $Y$ in the
decomposition, use the fixed smooth mother function $\eta$ on $[1,4]$
and $\eta_Y(t)=\eta(t/Y)$. Then $\log Y\gg L^\alpha$ and
$\log Y\ll_k L^\beta$. In particular $Y=N^{o(1)}$ and $\Delta Y=o(z)$.

For clarity, we give the precise normalization of the square that will
be estimated. With $q_0$ the integer cofactor, put
\[
 S_Y(q_0)=\sum_{\substack{a\in\mathcal A,\ m,n\ \text{satisfy }\eqref{eq:physblock}\\
                          mn-\Delta Q=ar_1q_0}}
       \eta_Y(a)\alpha_m\beta_n,
 \qquad \mathcal R_Y=\mathcal V^{-2}\sum_{q_0\in\mathbb Z}|S_Y(q_0)|^2.
\]
The original sum has $q_0>0$ and a bounded cofactor coefficient. There
are $O(X/Y)$ possible positive cofactors, so its square is bounded by
$O((X/Y)\mathcal R_Y)$. Terms of the expanded square satisfy
$a q_0\equiv b q_0\equiv Q\pmod{r_2}$, hence $a\equiv b\pmod{r_2}$.
When $(a,b)=1$, using the numerator divisibility by $r_1$ from
\eqref{eq:physblock}, the common integer cofactor condition is equivalent to
$bmn-ars=\Delta Q(b-a)$. The implication from the determinant identity
uses $(a,b)=1$ to deduce $a\mid(mn-\Delta Q)/r_1$ and
$b\mid(rs-\Delta Q)/r_1$.

Put $Z=\Delta Y$ and $B=L^{A_0}$. On the circle $\mathbb T=\mathbb R/\mathbb Z$
with mass one let $\mathfrak M$ be the union of the arcs
$\|\theta-h/q\|\le2B/Z$ over reduced $h/q$ with $q\le B$, and
$\mathfrak m=\mathbb T\setminus\mathfrak M$. Distinct such rationals
have spacing at least $B^{-2}$, so the arcs are disjoint for large $N$
because $B^3=o(Z)$. Choose a fixed smooth even
$\psi$, equal to one on $[-3,3]$ and supported on $[-4,4]$, with values in
$[0,1]$. Write $e(t)=\exp(2\pi i t)$ and, for $\mathcal I=\mathbb T,\mathfrak M,\mathfrak m$, set
\begin{equation}\label{eq:arckernel}
 K_{\mathcal I}(u;a,b)=\psi(u/Z)
    \int_{\mathcal I} e\bigl(\theta[u-\Delta(b-a)]\bigr)\,d\theta.
\end{equation}
Use this kernel only if $Q\mid t=bmn-ars$, taking $u=t/Q$.
Write $K_{\mathrm{raw}},K_{\mathrm{maj}},K_{\mathrm{min}}$ for the three
kernels. They are supported on $|u|\le4Z$, $|K_{\mathrm{maj}}|\ll B^3/Z$,
and $K_{\mathrm{raw}}=\mathbf1_{u=\Delta(b-a)}$ on the label support.
Define $\mathcal M_Y$ as the normalized sixfold sum with kernel
$K_{\mathrm{maj}}$ and the projections \eqref{eq:physblock} and
$a\equiv b\pmod{r_2}$. Define $\mathcal N_Y=\mathcal R_Y-\mathcal M_Y$.
Shared-label terms of $\mathcal R_Y$ always use the common integer
cofactor definition.

Here are the absolute estimates that justify subsequent deletions.
For fixed $a,b\asymp Y$, every row and column in the product-pair
support, with products $K_0=mn$ and $K_1=rs$, has $O(1+X/Q)=O(R)$
entries. Indeed, $Q\mid bK_0-aK_1$ fixes a class modulo $Q$ in either
variable, because $a,b$ are units, and the variables are at most $4X$.
Thus the weighted count with divisor weights is $O(RXL^3)$, by
Cauchy--Schwarz and $\sum_{n\le x}\tau(n)^2\ll x\log^3(2x)$. If a
fixed active prime $\ell$ divides either product the bound improves
to $O(RXL^3\ell^{-1/2})$, by $\tau(\ell h)\le2\tau(h)$ and
Cauchy--Schwarz on the two marginals.

For a common cofactor and $K=ar_1q_0+\Delta Q$ in the product box,
$(ar_1,\Delta Q)=1$. The divisor inequalities
\[
 \tau(K)^2\le\tau_4(K)\le
 4\sum_{e\mid K,\ e\le (4X)^{3/4}}\tau_3(e)
\]
follow by taking the largest factor in each ordered quadruple.
Each $e$ with any solution specifies one class of $q_0$ modulo $e$.
The elementary estimates for the partial sums of $\tau_3(e)$ and
$\tau_3(e)/e$ give
\begin{equation}\label{eq:cofactorsecondmoment}
 \sum_{q_0}\tau(ar_1q_0+\Delta Q)^2
 \ll (X/Y)L^3+X^{3/4}L^2\ll (X/Y)L^3,
\end{equation}
where the sum is over the cofactor range in the box. For $\ell\nmid ar_1$
and $\ell\mid K$, first restrict $q_0$ to its class modulo $\ell$ and
apply the same argument to $K/\ell$. Its progression has step $ar_1$
and a coprime intercept. This gives $O((X/(Y\ell))L^3)$, since
$Y,\ell=N^{o(1)}$ and $X\ge N^{1-o(1)}$.

There are $O(Y^2\exp(-\lambda_1))$ pairs of labels that share an
active prime, by summing $\lfloor4Y/\ell\rfloor^2$ over such primes.
Estimate their raw terms by \eqref{eq:cofactorsecondmoment} and
Cauchy--Schwarz, and their major-arc terms by the product-pair bound
and $B^3/Z$. Both cost $O_A(XYL^{-A})$ for every fixed $A$.
For coprime labels the raw identity rules out $\ell\mid(b,s)$ and
$\ell\mid(a,n)$, since reducing modulo $\ell$ would give
$\ell\mid\Delta Q$ contrary to the active-prime support.
On the major arcs the corresponding deletions cost $O_A(XYL^{-A})$
by the product-pair bound with $\ell^{-1/2}$, summing over at most $2k$
primes for each pair. When $t=0$, coprimality and roughness imply
$b\mid s$, so the remaining terms exclude $t=0$. These arguments also
give absolute bounds $O(XYL^{4C+3})$ and $O(RXYL^{4C+3}B^3)$ for the raw
and major sums, respectively.

Independently in each active interval choose $J=\lfloor L^\jmath\rfloor$
ordered auxiliary primes, with replacement and probability
$(\ell V_j)^{-1}$ per slot. Let $d$ be their product and let $D$ be the
unique dyadic number with $D\le d<2D$. A sample has weight $1/(dV^J)$.
Keep only squarefree $d$ coprime to $abns$. A union bound over slot
pairs and deleted label values costs $O_k((J+1)^2e^{-\lambda_1})$ of
the mass. For a fixed auxiliary prime $\ell\nmid ab$, the terms with
$\ell\mid n$ or $\ell\mid s$ gain $\ell^{-1/2}$ by the modified
\eqref{eq:cofactorsecondmoment} and the modified product-pair bound.
After averaging each slot against its probability and summing the
$kJ$ slots, these costs are again $O_A(XYL^{-A})$. The later moment retains the signed kernels.

\subsection{The operator with its exact measure}

Let $A_D^{\rm box}=DYM$, $\mathcal B_D=[A_D^{\rm box},16A_D^{\rm box}]\times[B_0,2B_0]$,
and $\kappa_J=((J+1)!)^k \mathcal V^{-J-1}$. A physical state consists of a
primitive point $\mathbf v=(x,y)^{\mathsf T}\in\mathcal B_D\cap\mathbb Z^2$
whose coordinates are units modulo $Q$, and a product $c\mid x$ of
$J+1$ distinct primes in each active interval. Require
\begin{equation}\label{eq:stateprojections}
 x\equiv ic,\qquad y\equiv r_2Qi^{-1}\pmod{r_1},
 \qquad y\equiv0\pmod{r_2}.
\end{equation}
Give every state measure $\kappa_J$ and put
$f_D(\mathbf v,c)=\overline{\alpha_{x/c}}\beta_y$ with zero extension of
the box coefficients. At a point where a value is nonzero, roughness
forces $c$ to contain precisely the active primes of $x$, each once.
Thus $f_D$ is constant on the state labels of each physical point,
and we denote this value by $f_D^\circ(\mathbf v)$.

The transition $T_D$ chooses $b\in\mathcal A$ dividing $c$ with weight
$(J+1)^{-k}$ and sets $d=c/b\in[D,2D)$. Sum with weight $\mathcal V^{-1}$ over
$a\in\mathcal A$ with $(a,c)=1$, and over target states
$(\mathbf v',ad)$ satisfying $Q\mid\det(\mathbf v,\mathbf v')$ and
$c\equiv ad\pmod{r_2}$. The edge carries the remaining factor
\begin{equation}\label{eq:physicaledgeweight}
 \frac Dd\,\eta_Y(a)\eta_Y(b)
 K_{\mathrm{min}}\!\left(\frac{\det(\mathbf v,\mathbf v')}{dQ};a,b\right).
\end{equation}
The quotient is integral since $d$ divides both first coordinates
and is coprime to $Q$. The congruence is equivalent to
$a\equiv b\pmod{r_2}$. Reversing the edge swaps $a,b$, and the kernel
satisfies $K_{\mathrm{min}}(-u;b,a)=\overline{K_{\mathrm{min}}(u;a,b)}$.
The constant state measure thus makes $T_D$ self-adjoint.

The physical correspondence is
\begin{equation}\label{eq:physicalvectors}
 \begin{gathered}
 \mathbf v=(bdm,s)^{\mathsf T},\quad \mathbf v'=(adr,n)^{\mathsf T},
 \quad c=bd,\quad c'=ad,\\
 \det(\mathbf v,\mathbf v')=d(bmn-ars).
 \end{gathered}
\end{equation}
The coordinate units follow from the coefficient masks and the label
support. A prime common to both coordinates of an endpoint would have
to divide its rough factor, by the auxiliary and endpoint deletions.
It would therefore be greater than $z$ and coprime to $Qd$, yet divide
$u=t/Q$ with $0<|u|\le4Z=o(z)$. This proves physical primitivity.
Conversely the edge recovers $m=x/c,r=x'/c',s=y,n=y'$.
Its endpoint primitivity supplies $(b,s)=(a,n)=1$ and $(d,ns)=1$,
and the label conditions supply $(a,b)=(d,ab)=1$.

With the convention $\langle f,g\rangle=\kappa_J\sum\overline f g$,
the coefficient product is exactly
$\alpha_m\overline{\alpha_r}\beta_n\overline{\beta_s}$. The common
weight in $D^{-1}\langle f_D,T_Df_D\rangle$ is
\begin{equation}\label{eq:pairingweight}
 \frac1D\frac{((J+1)!)^k}{\mathcal V^{J+1}}\frac1{(J+1)^k}\frac1{\mathcal V}\frac Dd
   =\frac{(J!)^k}{dV^{J+2}}.
\end{equation}
There are $(J!)^k$ ordered samples representing any admissible $d$.
The edge sum thus reproduces each surviving small-arc summand with the
probability weight of its admissible samples. In particular,
\begin{equation}\label{eq:operatorpairing}
 \mathcal N_Y=\sum_D D^{-1}\langle f_D,T_Df_D\rangle+O_A(XYL^{-A}).
\end{equation}
We will use the congruence projection on transitions repeatedly. For
orthogonal source and target projections $P_\gamma,P'_\gamma$ indexed
by $c\bmod r_2$, the operator
$\sum_\gamma P'_\gamma K P_\gamma$ has norm at most $\|K\|$.
Indeed its squared norm on a vector is the sum over $\gamma$ of the
squared norms of the separate outputs, each bounded by
$\|K\|^2\|P_\gamma f\|^2$. This proves that the edge matching can be
retained in all the norm bounds below.

\section{The major arcs}\label{sec:major-arcs}

\begin{proposition}\label{prop:major-arcs}
Assume Dirichlet GRH. For every fixed $E>0$, the projected major-arc
form satisfies $|\mathcal M_Y|\ll_E XYL^{-E}$, uniformly in the
reduced moduli, prime pairs, and Type~II boxes under consideration.
\end{proposition}

\begin{proof}
The local projections of the physical form remain in force. The first
coefficient has the form
\[
 \alpha_m=\mu(m)R_z(m)m^{it_0}w_m(m)\mathbf1_{(m,N)=1},
 \qquad |t_0|\leq N^{O(1)},
\]
where $w_m$ has uniformly bounded supremum and total variation.
Conditions on $m$ modulo $r_1$ can be included by orthogonality on
the units modulo $r_1$. The second coefficient has size $O(L^C)$ and
carries its own restrictions. The mean-square bounds below allow
both coefficients to have size $O(L^C)$.

Write $P_1=bmn$, $P_2=ars$, and $t=P_1-P_2$. Fix a reduced fraction
$h/q$ with $q\leq B$, and write $\theta=h/q+\xi$ with
$|\xi|\leq 2B/Z$. Factor $q=q_\parallel q_\perp$, where $q_\parallel$
contains the full prime powers in $q$ based at primes dividing $Q$.
Set $M_*=Qq_\parallel$. Both $P_1$ and $P_2$ are units modulo $M_*$.
For each $v\bmod q_\parallel$ and unit $\gamma\bmod q_\parallel$, impose
$P_2\equiv\gamma\pmod{q_\parallel}$. On that class,
\begin{equation}\label{eq:lift-congruence}
 t\equiv Qv\pmod{M_*}
 \quad\Longleftrightarrow\quad
 P_1P_2^{-1}\equiv 1+Qv\gamma^{-1}\pmod{M_*}.
\end{equation}
Indeed, $P_2^{-1}$ and any lift of $\gamma^{-1}$ agree modulo
$q_\parallel$. Multiplication by $Qv$ makes the right-hand residue
independent of this lift. It is a unit since every prime of $M_*$
divides $Q$. These classes partition $Q\mid t$. The indicator of
the right-hand congruence in \eqref{eq:lift-congruence} is
\begin{equation}\label{eq:lift-character}
 \frac{1}{\varphi(M_*)}\sum_{\chi\bmod M_*}
 \chi(P_1)\overline{\chi(P_2)}
 \,\overline{\chi(1+Qv\gamma^{-1})}.
\end{equation}
The case $q_\parallel=1$ has a single residue.

Put $e_c(u)=e(u/c)$. By the Chinese remainder theorem, on each class
the rational phase equals
\[
 e_{Qq}(ht)=e_{q_\parallel}(h\overline{q_\perp}v)
             e_{q_\perp}(h\overline{M_*}(P_1-P_2)).
\]
Each inverse is taken in the modulus of its factor. The second factor
separates at the endpoints. For example,
\begin{align*}
 e_{q_\perp}(h\overline{M_*}bmn)
  &=\sum_{\nu\bmod q_\perp}\nu(bm)\Omega_\nu(n),\\
 \Omega_\nu(n)&=\frac{1}{\varphi(q_\perp)}
  \sum_{u\bmod q_\perp}^{*}\overline{\nu(u)}
       e_{q_\perp}(h\overline{M_*}un).
\end{align*}
Here $bm$ is a unit modulo $q_\perp$, and $|\Omega_\nu(n)|\leq1$
for all $n$. The additive sign reverses at $ars$. The label phase
$e_q(-h\Delta(b-a))$ likewise separates on the label unit classes.
All active primes exceed $B$, and $m,r$ are $z$-rough, which gives
the required units. Expand $P_2\equiv\gamma\pmod{q_\parallel}$
multiplicatively. For $a\equiv b\pmod{r_2}$, sum over their common unit
residue and expand both indicators. The $n,s$ projections stay on
the second coefficient. For fixed $h/q$ the number of separated
terms is at most $B^{10}R^8$. The pair $v,\gamma$ uses at most
$q_\parallel^2$ choices, the projection at $P_2$ at most $q_\parallel$,
the product phases at most $q_\perp^2$, the label phases at most $q^2$,
and the local prime classes at most $r_2^3r_1^2$. All small characters
have moduli dividing $qr_1r_2$ and the expansion coefficients have
modulus at most $1$, with the normalization in \eqref{eq:lift-character}.

To separate the continuous weight, set $S_0=XY$, $s_j=\log(P_j/S_0)$,
and $\lambda=X/(Q\Delta)\ll R$. Both $P_j/S_0$ lie in $[1,16]$.
Choose fixed compactly supported smooth functions $W_1,W_2$ equal to
one on $[0,\log 16]$. The coupled factor is
\[
 F(s_1,s_2)=W_1(s_1)W_2(s_2)
   \psi\!\left(\lambda(e^{s_1}-e^{s_2})\right)
   e\!\left(\xi Z\lambda(e^{s_1}-e^{s_2})\right).
\]
With $U_*=(1+B)(1+\lambda)$, derivatives of order $j$ are
$O_j(U_*^j)$. Integration by parts gives $\|\widehat F\|_1\ll U_*^4$
and, for any integer $j>2$ and $T_0\geq1$,
\[
 \int_{\max(|u_1|,|u_2|)>T_0}|\widehat F(u_1,u_2)|\,du_1\,du_2
   \ll_j U_*^j T_0^{2-j}.
\]
At $T_0=L^{A_0+K+2}$ this tail has any fixed log-power saving by
choosing $j$ large. Fourier inversion produces the separated powers
$(P_1/S_0)^{2\pi i u_1}(P_2/S_0)^{2\pi i u_2}$. The slow phase
$e(-\xi\Delta(b-a))$ also separates into label weights and has
variation $O(B)$ on the range of either label.

Each term is now a character average of two label sums and two
product sums, with unit-modulus factors from \eqref{eq:lift-character}.
Denote a product sum by $U_\chi=\sum_{j\leq4X}u_j\chi(j)$, with
$|u_j|\ll L^{2C}\tau(j)$. Orthogonality and the bound $O(1+X/M_*)$
for each product residue fiber give
\begin{equation}\label{eq:major-product-moment}
 \frac1{\varphi(M_*)}\sum_{\chi\bmod M_*}|U_\chi|^2
 \ll (1+X/M_*)XL^{4C+3}\ll RXL^{4C+3}.
\end{equation}
This bound holds at every Fourier height. At each end the trivial
label bound is $O(Y)$. By Cauchy--Schwarz, the discarded Fourier tail
after all expansions and arc integration costs at most
\[
 XYL^{4C+13A_0+9K+3}
    \int_{\max(|u_1|,|u_2|)>T_0}|\widehat F(u_1,u_2)|\,du_1\,du_2,
\]
with a uniform bound for the last integral. We used $Z=\Delta Y\geq Y$,
total arc measure $O(B^3/Z)$, and $\mathcal V^{-2}=O_k(1)$.

Let $Q_s\mid qr_1r_2$ be the least common multiple of the small moduli.
If $\chi$ times the small label character at either endpoint is
principal on the common unit group, then $\chi$ is induced from
$(M_*,Q_s)$. On each prime-power component the kernel of reduction to
$Q_s$ from the units modulo $\operatorname{lcm}(M_*,Q_s)$ projects
onto the kernel of reduction from $M_*$ to $(M_*,Q_s)$. The product
being principal forces $\chi$ to be trivial on that last kernel.
Thus at most $\varphi((M_*,Q_s))\leq BR^2$ characters are exceptional.
This list includes the principal character.

For a nonexceptional character, fix the product $b'$ of $k-1$ label
factors. The last prime lies in its active interval and in
$[Y/b',4Y/b']$. Its character is nonprincipal and has modulus dividing
$\operatorname{lcm}(M_*,Q_s)\leq NL^{O(1)}$. The Dirichlet GRH
explicit formula gives $O(\sqrt{x}\log^2(2x\operatorname{lcm}(M_*,Q_s)))$
for the character-weighted von Mangoldt sum at endpoint $x$. Use
truncation height $x$ and the Dirichlet zero count to sum the inverse
zero ordinates. Passing to an inducing character changes the sum by
$O(\log(2\operatorname{lcm}(M_*,Q_s))\log(2x))$.
\cite[Lemma~2]{KandhilLanguascoMoree}
After removing higher prime powers and using Abel summation, the
prime sum is $O((Y/b')^{1/2}L^{O(1)})$. Here the fixed rescaling
$\eta_Y$ has bounded variation, and the slow phase and Fourier twist
have variation $L^{O(1)}$ for $|u_j|\leq T_0$. Removing the primes
dividing $N$ has additional cost $L^{O(1)}$. On a nonempty range
$Y/b'\gg e^{L^\alpha}$, and $\sum_{b'}1/b'\ll_k1$. Therefore either
label sum is $O_k(YL^{O(1)}e^{-cL^\alpha})$ for some $c>0$.
By \eqref{eq:major-product-moment}, the nonexceptional contribution
is $O_E(XY L^{-E})$ after all expansions and integrations.

For an exceptional $\chi$, the character in the first coefficient
equals a small-modulus character on its support $(m,N)=1$. Removing
that mask adds only $z$-rough integers divisible by a prime
$\ell\mid N$, $\ell>z$. Their count in $(M,2M]$ is
$O(LM/z+L)$, smaller than $ML^{-D_m}$ for every fixed $D_m$.
Apply \cite[Proposition~8.1]{Lichtman} to the resulting rough
M\"obius sum, with bounded-variation weight and twist height
$N^{O(1)}$. It is $O_{D_m}(ML^{-D_m})$. Applying this saving at one
endpoint and trivial estimates on the other factors gives the
exceptional contribution before small-modulus expansions and
integration as
\[
 \ll_{D_m} X^2Y^2 L^{4C-D_m}\frac{BR^2 L}{Q},
\]
using $1/\varphi(M_*)\ll L/Q$. The total cost after division by $XY$
is at most $L^{4C+18A_0+15K+1-D_m}$, using
$\|\widehat F\|_1\ll(1+B)^4(1+\lambda)^4$, the arc measure, and
$X/Q\ll R$. Take $D_m>E+4C+20A_0+16K+30$ and the Fourier-tail saving
larger than $E+4C+13A_0+9K+20$.
\end{proof}

\section{Unit-slope lattices and root equidistribution}\label{sec:root-count}

The roots used below are primitive integer vectors whose coordinates
are both units modulo $Q$.

For a physically primitive root $\mathbf v=(x,y)^{\mathsf T}$, solve
$yx_0\equiv Q\pmod{x}$ with $0\leq x_0<x$. Put $y_0=(yx_0-Q)/x$.
Then
\begin{equation}\label{eq:lattice-completion}
\begin{gathered}
 G=\begin{pmatrix}x&x_0\\y&y_0\end{pmatrix},\quad
 \det G=-Q,\quad \rho=x_0/x,\quad
 \tau_\rho(\mathbf w)=w_1+\rho w_2,\\
 G\mathbf w=\binom{x\tau_\rho(\mathbf w)}
                    {y\tau_\rho(\mathbf w)-Qw_2/x}.
\end{gathered}
\end{equation}
The image of $G$ is the index-$Q$ lattice
\[
 \Lambda_\lambda=\{(u,v)^{\mathsf T}\in\mathbb Z^2:
                   v\equiv\lambda u\pmod Q\},
 \qquad \lambda=yx^{-1}\pmod Q .
\]
Both columns have this slope, and $(u,v)\mapsto v-\lambda u\bmod Q$
is surjective, so its kernel and $G\mathbb Z^2$ have the same index.
For a point $G\mathbf w$ whose first coordinate is a unit modulo
$Q$, physical primitivity is equivalent to primitivity of $\mathbf w$.
A prime dividing both coordinates of $G\mathbf w$ is coprime to $Q$.
Multiplying by the adjugate shows that it divides both coordinates
of $\mathbf w$. The converse follows by matrix multiplication.

For physical vertices the condition $Q\mid\det(\mathbf v,\mathbf v')$
is equivalent to equality of the unit slopes. Their paths thus have
a common lattice.

Let $A_D^{\rm box}=DYM$ and $B_0$ be the scales of the two physical
coordinates. Then $A_D^{\rm box}B_0=DYX$, both scales are at least
$N^\delta$, and $4X\geq\Delta Q$. The root box is
$[A_D^{\rm box},16A_D^{\rm box}]\times[B_0,2B_0]$.
To exchange coordinates when $B_0<A_D^{\rm box}$, set
$\epsilon=Q/(xy)\leq4/(DY\Delta)<1$ and $j=\lceil\rho-\epsilon\rceil$.
With $P=\bigl(\begin{smallmatrix}0&1\\1&0\end{smallmatrix}\bigr)$ and
$R_j=\bigl(\begin{smallmatrix}1&j\\0&-1\end{smallmatrix}\bigr)$, the
new normalized matrix and parameter are
\begin{equation}\label{eq:coordinate-swap}
 G'=PGR_j,\qquad \rho'=j-\rho+\epsilon,\qquad G'R_j\mathbf w=PG\mathbf w .
\end{equation}
The parameter $\rho'$ lies in $[0,1)$. The branch is involutive
because $R_j^2=1$ and $\lceil\rho'-\epsilon\rceil=j$. Both coordinates
remain units, and for any modulus $T$ the matrix class transforms
by $C\mapsto PCR_j$. This preserves the determinant and sends the
old second physical row to the new first row. The absolute Jacobian
in normalized real coordinates $(x/A_D^{\rm box},y/B_0,\rho)$ is
$1$. The branch boundary $\rho=\epsilon$ has bounded slopes in these
coordinates.

\begin{proposition}\label{prop:root-equidistribution}
Let $T=N^{o(1)}$ be squarefree and coprime to $Q$. Fix a matrix class
$C\bmod T$ of determinant $-Q$. For a rectangle in the normalized
three coordinates, the number of primitive roots with $(xy,Q)=1$
and $G\equiv C\pmod T$, divided by $A_D^{\rm box}B_0$, is the volume
of the rectangle times the density
\begin{equation}\label{eq:root-density}
 \frac{c(Q)}{\zeta(2)|\mathrm{SL}_2(\mathbb Z/T\mathbb Z)|},\qquad
 c(Q)=\prod_{\ell\mid Q}\frac{\ell-1}{\ell+1}\leq1,
\end{equation}
with error $O(N^{-\delta/1000+o(1)})$ uniform in the rectangle and class.
\end{proposition}

\begin{proof}
First put the short side in the first coordinate, and write $A_s,B_l$
for the short and long scales. Fix $x\equiv C_{11}\pmod T$ coprime
to $Q$. Put $T_1=(x,T)$ and $T_2=T/T_1$. After imposing the classes
of $y,x_0$, the congruence prescribing $y_0\bmod T$ reduces to
\begin{equation}\label{eq:completion-congruence}
 yx_0\equiv A_x:=Q+xC_{22}\pmod{xT_1}.
\end{equation}
The congruence modulo $T_2$ follows from $\det C=-Q\bmod T_2$.
The integer $A_x$ is a unit modulo $xT_1$, whose prime divisors are
those of $x$.

Apply exact M\"obius inversion
$\mathbf1_{(y,Q)=1}=\sum_{d'\mid(Q,y)}\mu(d')$. For each $d'\mid Q$
write $y=y_*+d'Tj$ and $x_0=C_{12}+Tk$, where $y_*$ imposes
$d'\mid y$ and $y\equiv C_{21}\pmod T$. Because
$(d'T,xT_1)=(T,xT_1)=T_1$, both parameters have period $x$ in
\eqref{eq:completion-congruence}. The number of permitted pairs
over one period is
\begin{equation}\label{eq:zero-frequency-count}
 x\prod_{\ell\mid x,\,\ell\nmid T}(1-\ell^{-1}).
\end{equation}
Indeed, the pairs correspond to units $y\bmod xT_1$ with
$y\equiv C_{21}\pmod{T_1}$. Each gives $x_0=A_xy^{-1}$, which has
the required class modulo $T_1$. Here the determinant class ensures
$C_{21}$ and $C_{12}$ are units at primes dividing $T_1$. The unit
condition in \eqref{eq:completion-congruence} makes $(x,y)=1$ automatic.

Periodize the two parameter intervals modulo $x$. Their Fourier
coefficients, normalized by the period, are $O((1+|m|)^{-1})$ for
symmetric nonzero frequency representatives $|m|\leq x/2$.
The sum over any full period vanishes. The remaining consecutive
geometric sum divided by $x$ is bounded by
$x^{-1}\min(\text{length},2/|1-e(m/x)|)$. The zero coefficient is
the parameter interval length divided by $x$, with $O(x^{-1})$
rounding error. In the long coordinate it is $O(1+B_l/x)$ uniformly
in $d'$.

Apply an additive character projection to $y\equiv C_{21}\pmod{T_1}$.
For parameter frequencies $(m,n)$ the completed sums are
Kloosterman sums modulo $xT_1$ with frequencies
\[
 m\overline{d'T_2}+hx,\qquad n\overline{T_2}A_x,
\]
up to constant phases. The inverses here are modulo $xT_1$.
Reducing the frequencies modulo $x$ shows that their common divisor
with $xT_1$ is at most $T_1(m,n,x)$. For all moduli $c$ use the
Weil--Kloosterman bound
\[
 \left|\sum_{u\bmod c}^{*} e_c(mu+nu^{-1})\right|
 \leq\tau(c)c^{1/2}(m,n,c)^{1/2}.
\]
\cite[Eq.~(1.2)]{BlomerMan}
For $m,n\ne0$, grouping by divisors gives
\[
 \sum_{m,n\ne0}\frac{(m,n,x)^{1/2}}{(1+|m|)(1+|n|)}
 \ll \log^2(2x)\sum_{d\mid x}d^{-3/2}=x^{o(1)}
\]
over the symmetric representatives. With one frequency zero the
factor is $\log(2x)\sum_{d\mid x}d^{-1/2}=x^{o(1)}$, times the
zero Fourier coefficient of the other interval. The nonzero
frequency error for fixed $x,d'$ is thus
$O(x^{1/2+o(1)}T^{O(1)}(1+B_l/x))$. Summing over $x,d'$ with $\mu(d')\ne0$ and
dividing by $A_sB_l$ bounds the error by $A_s^{-1/2}N^{o(1)}$, since
$B_l\geq A_s$ and $2^{\omega(Q)}=N^{o(1)}$. This completion applies
also when the $d'$ parameter interval has length less than one.

The zero Fourier term in \eqref{eq:zero-frequency-count} contributes
the long-side interval length divided by $d'T$ and the normalized
$\rho$ interval length divided by $T$, times the product in that
formula. Rounding costs $A_s^{-1}N^{o(1)}$ after normalization.
Truncating the zero-frequency divisor main term at $d'\leq U$
has normalized error at most a constant times
\begin{equation}\label{eq:large-divisor-error}
 \sum_{\substack{d'\mid Q\\d'>U}}\frac{|\mu(d')|}{d'}
   \leq \frac{2^{\omega(Q)}}{U}.
\end{equation}
Take $U=N^{\delta/4}$, keeping the uniform completed error for all
$d'$ and restoring the full main term with the same bound. The
sieve factor on the long coordinate is therefore
$\sum_{d'\mid Q}\mu(d')/d'=\varphi(Q)/Q$.

On the short side expand $\prod_{\ell\mid x,\,\ell\nmid T}(1-\ell^{-1})$
as $\sum_{e\mid x,\,(e,T)=1}\mu(e)/e$. Since $(x,Q)=1$, it suffices to
sum over $(e,TQ)=1$. Expanding the condition $(x,Q)=1$ by M\"obius
inversion gives a progression count with density
$T^{-1}e^{-1}\varphi(Q)/Q$ for each $e$, and error $O(2^{\omega(Q)})$.
Summing the errors with their weights costs
$O(2^{\omega(Q)}\log(2A_s))$ before short-side normalization.
Extending the main sum in $e$ to infinity costs $O(A_s^{-1})$ in
relative density. The density in the three coordinates is now
\[
 T^{-3}\left(\frac{\varphi(Q)}{Q}\right)^2
       \prod_{\ell\nmid TQ}(1-\ell^{-2}).
\]
For squarefree $T$,
$|\mathrm{SL}_2(\mathbb Z/T\mathbb Z)|=T^3\prod_{\ell\mid T}(1-\ell^{-2})$.
This gives \eqref{eq:root-density}.

If the coordinates were exchanged, approximate the image of a
rectangle on each branch of \eqref{eq:coordinate-swap} by mesh boxes
of side $N^{-\delta/1000}$. The boundary consists of finitely many
graphs of bounded slope, so their boundary-box volume is
$O(N^{-\delta/1000})$. There are $O(N^{3\delta/1000})$ boxes.
The short-first counting and large-divisor errors remain within
$O(N^{-\delta/1000+o(1)})$ after summation. The Jacobian and matrix-class
bijection in \eqref{eq:coordinate-swap} give the same density in the
original coordinates.
\end{proof}

\section{Relative paths and the history expansion}\label{sec:history}

\subsection{The enlarged state space and good roots}

We use the lattices and root count of Section~\ref{sec:root-count}.
Put $\Lambda_\lambda=\{(x,y)\in\mathbb Z^2:y\equiv\lambda x\pmod Q\}$
for a unit slope $\lambda\bmod Q$. A physical state belongs to exactly
one such lattice and is primitive in it. In each lattice take a copy
of every state satisfying the same box and label conditions, with
primitivity in that lattice in place of physical primitivity and no
coordinate unit requirement. Edges are within one copy. We extend
$f_D$ by zero off the physical states. On physical states, determinant
divisibility is equivalent to equality of unit slopes. Therefore
\eqref{eq:operatorpairing} is unchanged by the enlargement, and
self-adjointness continues to hold.

For a physical root $(x,y)$ let $G$ be the completion of determinant
$-Q$ in \eqref{eq:lattice-completion}, with first column $(x,y)^{\mathsf T}$.
Write $\rho=x_0/x$ and $\tau_\rho(\mathbf w)=w_1+\rho w_2$. In the
same lattice every point has a unique relative coordinate $\mathbf w$
under $G$, and lattice primitivity means $\gcd(w_1,w_2)=1$.
On an edge the kernel argument in \eqref{eq:physicaledgeweight} becomes
$-\det(\mathbf w,\mathbf w')/d$. If
$\det(\mathbf w,\widehat{\mathbf w})=1$, the normalized completion
parameter at $G\mathbf w$ is
\begin{equation}\label{eq:rhotransfer}
 \frac{\tau_\rho(\widehat{\mathbf w})}{\tau_\rho(\mathbf w)}\pmod1.
\end{equation}
Indeed $G\widehat{\mathbf w}$ completes the new point with determinant
$-Q$, and adding an integer multiple of the new point normalizes its
first coordinate. This also defines the parameter at every state of
the enlarged space. Geometric projections refer to both coordinates
of $G\mathbf w$, including $-Qw_2/x$ in the second coordinate. The
local projections are \eqref{eq:stateprojections} at that point. The
root has $y\equiv0\pmod{r_2}$ and $xy_0\equiv-Q\pmod{r_2}$, so the
$r_2$ condition on the second coordinate becomes $w_2\equiv0\pmod{r_2}$.

For a state with label $c$ and completion parameter $\rho_*$, impose
the following two conditions.
\begin{enumerate}
\item For all admissible $d=c/b\in[D,2D)$ and integers
 $1\le h\le Y^{1/5}$, we have $\|hd\rho_*\|>Y^{-7/10}$.
\item Let $I$ be any nonempty subset of the interval indices and
 $i_*=\max I$. For each $i\in I\setminus\{i_*\}$ choose a prime in
 $\mathcal P_i$ and denote their product by $P_<$. Independently
 draw one prime for each $i\notin I$ with weight $(\ell V_i)^{-1}$,
 and denote their product by $T_f$. Apart from a set of $T_f$ of
 draw mass at most $e^{-\lambda_{i_*}/4}$, the values
 $\Delta dP_<T_f\rho_*\pmod1$ for distinct candidates $dP_<$ are
 pairwise separated by at least $100e^{-\lambda_{i_*}}$.
\end{enumerate}
There are at most $(J+1)^k$ candidates $d$. The first test fails on a
set of measure at most $2(J+1)^kY^{-1/2}$. For each $I$ the separation
of the intervals bounds the number of $dP_<$ by
$(J+1)^k\exp(2\lambda_{i_*}/7)$. For fixed $T_f$, each colliding pair
has a nonzero integer difference multiplier, and its collision set
has Haar measure at most $200e^{-\lambda_{i_*}}$. A union bound and
Markov's inequality for the $T_f$ mass show, uniformly in $c$, that the
failed second tests together have measure $O_k(e^{-c_1L^\alpha})$
for a fixed $c_1>0$. This bound holds for both tests jointly.
Let $\mathsf G$ project onto the good states and set
$A_D=\mathsf G T_D\mathsf G$, which is self-adjoint.

\subsection{Using the root count for path configurations}

Set $H=\lceil L^\varrho/2\rceil$. A relative path of at most $2H$
edges starts at $\mathbf e_1=(1,0)^{\mathsf T}$. Its edge determinant
in physical coordinates is $O(QDZ)$. Since all first coordinates
are comparable to $A_D^{\rm box}$, the identity
\[
 \frac{y_{j+1}}{x_{j+1}}-\frac{y_j}{x_j}
 =\frac{\det(\mathbf v_j,\mathbf v_{j+1})}{x_jx_{j+1}}
\]
gives $|w_{j,2}|=|\det(\mathbf v_0,\mathbf v_j)|/Q\ll HDZ$.
Moreover $1/16\le\tau_\rho(\mathbf w_j)\le16$, so
$|\mathbf w_j|\ll HDZ$. For each configuration fix its labels,
relative vertices, and matrix class modulo the squarefree
$T=r_1r_2\prod_\ell\ell$, where the product is over its distinct
active primes. Then $(T,Q)=1$, the class fixes the divisibility and
local projections, and the kernels depend only on relative
determinants and labels. If $W_1$ is the total absolute configuration
weight plus one, then
\begin{equation}\label{eq:pathbudget}
 \begin{split}
 \log T+\log W_1 &\ll_k H\bigl(\log(HDZ)+(J+1)L^\beta \\
     &\hspace{4em} +(J+1)\log(J+1)+\log R\bigr)=o(L).
 \end{split}
\end{equation}
For this bound one may independently allow $O((HDZ)^2)$ choices for
each relative vertex and $\exp(O(L^\beta))$ choices for each prime
slot, and bound the remaining coefficients, factorials, and arc
kernels absolutely. It uses $\varrho+\beta+\jmath<1$.

To pass from rectangles in Proposition~\ref{prop:root-equidistribution}
to the required configuration indicators, take a mesh of side
$N^{-\delta/8000}$ in $(x/A_D^{\rm box},y/B_0,\rho)$. The position
boundary in the first or second normalized coordinate has derivative
$\tau_\rho(\mathbf w)$ in that coordinate, within the bounds above on
the relevant support. For the unwrapped parameter in
\eqref{eq:rhotransfer}, differentiation in $\rho$ gives
$\tau_\rho(\mathbf w)^{-2}$. An integer completion of size
$O(|\mathbf w|+1)$ can be chosen by B\'ezout's identity. All
integer multipliers and shifts occurring in the good-point tests,
and all possible prime choices in those tests, cut the boundaries
into $N^{o(1)}$ pieces of slope $N^{o(1)}$. The weighted exception
condition in the second test is constant between the individual
separation boundaries. The exchange boundary in
\eqref{eq:coordinate-swap} obeys the same bound. The total volume of
boundary mesh boxes is $N^{-\delta/8000+o(1)}$. Summing the rectangular
error over $O(N^{3\delta/8000})$ mesh boxes, then over configurations
and classes using \eqref{eq:pathbudget}, gives total normalized error
$N^{-\delta/8000+o(1)}$. Since $H\ll L^\varrho$ with $\varrho<1$, this
is $O_A(L^{-2AH})$ for every fixed $A$. The same calculation applies
to a single edge with no endpoint return condition.

Only physical roots are counted. Their extra density is the common
factor $c(Q)=\prod_{\ell\mid Q}(\ell-1)/(\ell+1)$. The CRT components
on active primes and at $r_1,r_2$ are independent. For each active
prime $\ell$, the first row of $G\bmod\ell$ has a uniformly random
kernel line. For each nonzero first row there are exactly $\ell$
second rows giving determinant $-Q$, and every line is the kernel of
exactly $\ell-1$ nonzero rows. Relative points are primitive and
have nonzero reduction. Thus the factor at $\ell$ is exactly
\begin{equation}\label{eq:linefactor}
 \frac1{\ell+1}\sum_{\mathcal L\in\mathbb P^1(\mathbb F_\ell)}
   \prod_{j:\ell\mid c_j}\mathbf1_{\overline{\mathbf w_j}\in\mathcal L}.
\end{equation}

For the removal of bad physical endpoints, put $(1-\mathsf G)f_D$
at either endpoint of the difference of the pairings for $T_D$ and
$A_D$, and use self-adjointness to treat the bad endpoint as the
root. Bound the endpoint coefficient product by $O(L^{4C})$. In a
one-edge configuration the raw kernel fixes $w_2=\Delta d(a-b)$,
while the major kernel allows $O(Z)$ values of $w_2/d$ at cost
$O(B^3/Z)$ each. There are $O(1)$ possible $w_1$ for each value.
The absolute target mass at fixed labels is $O(1+B^3)$.
In the uniform class average, the line factors give at most
$1/(\ell+1)$ for each distinct chosen prime. The total mass, after
dropping position and local restrictions, is at most
\[
 \mathcal V^{-(J+2)}\prod_{j=1}^k
  \left(\sum_{\ell\in\mathcal P_j}\frac1{\ell+1}\right)^{J+2}\le 1,
\]
by \eqref{eq:pairingweight}, bounding $D/d\le1$ in the undivided
pairing and reordering the auxiliary samples. The failure measure
is $O_k(e^{-c_1L^\alpha})$ at each fixed label. Thus the absolute
bad-root pairing divided by $DYX$ is at most
$O_k(L^{4C}(1+B^3)e^{-c_1L^\alpha})+O(N^{-\delta/8000+o(1)})$.
The dyadic $D$ values number at most $L$ for large $N$, since
$\log d\ll_k L^{\beta+\jmath}$. Consequently
\begin{equation}\label{eq:goodoperatorpairing}
 \mathcal N_Y=\sum_DD^{-1}\langle f_D,A_Df_D\rangle+O_A(XYL^{-A}).
\end{equation}

\subsection{The path moment and empty memory}

For a physical point $\mathbf v$, let $u_{\mathbf v}$ be the indicator
of all its state labels in its unit-slope copy. Consider
\[
 \mathcal T_H=\sum_{\mathbf v\,\text{physical}}
     \langle u_{\mathbf v},A_D^{2H}u_{\mathbf v}\rangle
 =\sum_{\mathbf v\,\text{physical}}\|A_D^Hu_{\mathbf v}\|^2.
\]
The root term is $\kappa_J$ times the sum over independent initial
and final labels of the matrix entry of $A_D^{2H}$ at the common
point. Expand matrix multiplication and transfer the edges to
relative coordinates. All paths now begin and end at $\mathbf e_1$.
The configuration average gives
\begin{equation}\label{eq:momentmean}
 \begin{aligned}
 \frac{\mathcal T_H}{DYX}&=\frac{c(Q)}{\zeta(2)}
  \int_1^{16}\int_1^2\int_0^1
     \mathbb E_{C_{r_1},C_{r_2}}\!\left[\mathcal M_H(A_D^{\rm box}\sigma,B_0\upsilon,\rho;C_{r_1},C_{r_2})\right]\\[-1pt]
     &\hspace{5em}\times d\rho\,d\upsilon\,d\sigma+O_A(L^{-2AH}).
 \end{aligned}
\end{equation}
Here $\mathcal M_H$ is the relative closed-path sum including the
factors \eqref{eq:linefactor}. The expectation is uniform on
determinant $-Q$ classes at $r_1,r_2$. The good-point, position,
and local projections remain in this sum.

Fix the real parameters and local classes in \eqref{eq:momentmean}.
Put $\nu_j(\ell)=((\ell+1)V_j)^{-1}$ and
$a_*=\max_{j,\ell}\nu_j(\ell)\ll_k e^{-\lambda_1}$, so
$\sum_{\ell\in\mathcal P_j}\nu_j(\ell)\le1$. A history state records
its relative point, active label $c$, and, for each interval, an
unordered set of temporarily absent pairs $(\ell,\mathcal L_\ell)$.
All active and temporarily absent primes are distinct at each state.
Cap the temporarily absent count at $M_0=\lceil L^2\rceil$.
The measure is the active weight
$\prod_j((J+1)!\prod_{\ell\mid c,\,\ell\in\mathcal P_j}\nu_j(\ell))$
times $\nu_j(\ell)$ per temporarily absent pair, with counting
measure on lines. An ordered list of $m_j$ absent pairs in interval
$j$ equivalently has weight divided by $m_j!$.

At each step choose the set $S$ of departing intervals whose prime
and current line are saved, and the set $I$ of incoming intervals
using a saved prime. A departure elsewhere is discarded. A return
comes from source memory, carries $V_i^{-1}$, and its line is checked
at the target. A fresh prime carries $\nu_i(\ell)$. An entry saved
at this step cannot return at that same step. Shared primes of $d$
require the same line at source and target. Keep the factor
$(J+1)^{-k}$, the weights and kernel of \eqref{eq:physicaledgeweight},
distinctness at an edge, and the geometric, good-state, and
congruence projections. Call the resulting operator $\mathsf E_j$.
Let $b_0$ indicate position $\mathbf e_1$ and empty memory, with
arbitrary active label. Its squared norm is at most
$\prod_j(\sum_{\ell\in\mathcal P_j}\nu_j(\ell))^{J+1}\le1$.

For a path in $\mathcal M_H$, a departing prime is saved exactly
when it later returns. The memory is empty at both endpoints, and
every prime has exactly one birth, including the initial entries.
Conversely an empty-to-empty history with distinct births supplies
this pattern of departures. Shared primes check the line over
consecutive active times, and saved pairs check it on each return.
These checks implement \eqref{eq:linefactor}. A prime with $h_\ell$
active stretches has weight $\nu_i(\ell)V_i^{1-h_\ell}$, equal to
$V_i^{-h_\ell}/(\ell+1)$. The initial factorials match $\kappa_J$,
and the edges have the original removal and $D/d$ weights. The
capacity is immaterial since $2kH<M_0$. Thus the path expansions
give the exact identity
\begin{equation}\label{eq:emptymemory}
 \mathcal M_H=
  \langle b_0,\mathsf E_0\mathsf E_1\cdots\mathsf E_{2H-1}b_0\rangle
    _{\text{all birth values distinct}}.
\end{equation}

\section{One-step history bounds}\label{sec:onestep}

At a primitive relative point $\mathbf w$, choose an integer
completion $\widehat{\mathbf w}$ with determinant one and write the
target as $j_0\mathbf w+j_1\widehat{\mathbf w}$. Its first-coordinate
ratio to the source is between $1/16$ and $16$, so each fixed $j_1$
allows $O(1)$ values of $j_0$. For fixed choices of primes, the raw
kernel fixes $j_1$ and the major kernel has $O(Z)$ possibilities
for $j_1/d$ at cost $O(B^3/Z)$ each. The total absolute target mass
is therefore $O(1+B^3)$. The returned slots have at most $M_0^k$
choices, birth weights have total mass at most one per interval,
and the departure label has normalized mass at most one. This bounds
the row sum by $O_k(M_0^k(1+B^3))$. Specifying $s$ fresh-slot values
gains the factor $a_*^s$ in this row bound.

For the weighted column bound reverse the edge. Source measure
times the forward birth weights equals target measure times the
reverse birth weights: a surviving prime keeps its measure, and
the active factorials are unchanged. The ratio of return weights
to reverse return weights is $\prod_{i\in S}V_i/\prod_{i\in I}V_i$,
which is bounded above and below in terms of $k$. The line checks
refer to the same occurrences in both directions. This proves the
same coarse bound for columns. In particular for
$c_*=2k+3A_0+10$, all absolute row and weighted column sums and the
Schur bound are at most $L^{c_*}$ for large $N$. The bound remains
valid with birth indicators and phases of modulus one attached.

\subsection{Steps with empty return and storage sets}

First allow an incoming and departing prime of the same interval
to coincide. Partition points into strips in $w_2/\tau_\rho(\mathbf w)$
of width $DZ$. The determinant support and the first-coordinate
bounds allow only $O(1)$ neighboring strips for each strip in
either direction. Fix $d$, the unchanged memory, a pair of such
strips and the lines at every prime dividing $d$. The common-line
lattice has index $d$. For a contributing pair choose a good
primitive reference point $\mathbf w_*$ in one of the strips.
With $\det(\mathbf w_*,\widehat{\mathbf w}_*)=1$, the lattice is
$h\mathbf w_*+jd\widehat{\mathbf w}_*$. Indeed the determinant with
$\mathbf w_*$ vanishes modulo every prime dividing $d$ exactly when
the residue lies on the prescribed line. The identity
\[
 \frac{w_2}{\tau_\rho(\mathbf w)}-
 \frac{w_{*,2}}{\tau_\rho(\mathbf w_*)}
 =\frac{\det(\mathbf w_*,\mathbf w)}%
         {\tau_\rho(\mathbf w_*)\tau_\rho(\mathbf w)}
\]
yields $|j|\ll Z$ and $|h+jd\rho_*|\ll1$, where
$\rho_*=\tau_\rho(\widehat{\mathbf w}_*)/\tau_\rho(\mathbf w_*)$.
The planar lattice with coordinates $(j/Z,h+jd\rho_*)$ has area
$1/Z$ per fundamental parallelogram. Its shortest vector has
length $\gg\Delta^{-1/5}Z^{-4/5}$ by the first good-point test:
use $|j|/Z$ when $|j|>Y^{1/5}$, the bound $Y^{-7/10}$ otherwise
for $j\ne0$, and $|h|\ge1$ when $j=0$. A reduced planar basis has
angle with sine bounded below and product of lengths $\asymp Z^{-1}$.
Consequently the two sets lie in a common coordinate rectangle of
sidelengths $A',B'$ with $A'B'\asymp Z$ and
$Z^{1/10}\le A',B'\le Z^{9/10}$. There is fixed exponent room here
since $\Delta\le L^K$ and $\log Y\gg L^\alpha$. Use the positively
oriented corresponding basis in the original lattice, whose
determinant is $d$. The relative determinant divided by $d$ becomes
$\det(\mathbf k,\mathbf l)$ in these coordinates.

With counting measure on this block, the two nonshared label factors
are $\eta_Y(a)\sqrt{m(a)}$ and $\eta_Y(b)\sqrt{m(b)}$, where
$m(a)=\prod_{\ell\mid a}\nu_i(\ell)\ll_k Y^{-1}$ on support, using the
interval $i$ of $\ell$ in each factor. This follows from the square-root
ratio of source and target measures times the forward birth weights.
All geometric, good-state, local vertex, and membership restrictions
are endpoint multipliers. The matching condition is a contraction by
the projection argument after \eqref{eq:operatorpairing}.
Substitute \eqref{eq:arckernel} and reverse the arc parameter on the
symmetric small arcs. Its phase separates into
$e(\theta\det(\mathbf k,\mathbf l))$ and $e(-\theta\Delta(a-b))$.
For any input $F$, extension by zero of the label frequencies and
Parseval give
\begin{equation}\label{eq:labelparseval}
 \int_{\mathbb T}\sum_{\mathbf k}
 \left|\sum_b \eta_Y(b)\sqrt{m(b)}F(\mathbf k,b)
             e(\pm\Delta b\theta)\right|^2\,d\theta
 \ll Y^{-1}\|F\|^2.
\end{equation}
The same holds at the other endpoint. Fourier-invert the smooth
cutoff $\psi(\det/Z)$ in a variable $v$ with weight $\widehat\psi(v)$.
This changes the geometric frequency to $\xi=\theta+v/Z$ while
keeping the label frequencies at $\theta$. By Cauchy--Schwarz and
\eqref{eq:labelparseval}, the norm at fixed $v$ is bounded by
$O(Y^{-1}\sup_{\theta\in\mathfrak m}\|\mathsf U_{\theta+v/Z}\|)$,
where $\mathsf U_\xi(\mathbf k,\mathbf l)=e(\xi\det(\mathbf k,\mathbf l))$.

If $|v|\le B$, distance from every reduced rational of denominator
$q\le B$ remains greater than or equal to $B/Z$, with strictness
from the small arcs. Dirichlet's approximation theorem at height
$\lceil Z/B\rceil$ supplies a reduced $h/q$ with
$|\xi-h/q|\le\min(B/Z,q^{-2})$ and $B<q\ll Z/B$.
On the common rectangle the geometric matrix is, up to exchanging
the input tensor factors, the tensor product of
$(e(\xi k_1l_2))$ and the transpose conjugate matrix. Its norm is
the squared norm of a one-dimensional rectangular Fourier matrix.
The Gram row sum is bounded by
$\sum_{|h'|\ll A'}\min(B',O(\|h'\xi\|^{-1}))$, interpreting the
summand as $B'$ at an integral argument. In each block of at most
$q$ consecutive $h'$, subtract the phase of the first index. The
approximant $h/q$ permutes the $1/q$ grid and the perturbation in
that block is at most $1/q$. The block sum is therefore
$O(B'+q\log(2q))$. Dividing the resulting norm by $Z$ yields
\begin{equation}\label{eq:geometryfourier}
 Z^{-1}\|\mathsf U_\xi\|\ll
 q^{-1}+Z^{-1/10}\log(2Z)+q\log(2q)/Z.
\end{equation}
The normalization needed in the step is $Y^{-1}$, giving the extra
factor $Z/Y=\Delta$. The $|v|\le B$ part is thus bounded by
$O(L^{K+\beta-A_0}+R Z^{-1/10}\log(2Z))$. For $|v|>B$ the trivial
matrix bound is $O(Z)$. The Schwartz tail gives $O_j(RB^{1-j})$
for each fixed $j>1$.

Restore the edge distinctness by subtracting the edges where a fresh
prime equals the departing prime. The target already excludes a
collision with $d$. Fixing the offending birth gives row and column
bounds $O_k((1+B^3)a_*)$ by the preceding counting and reversal.
The sum over $d$ uses the contraction
$F\mapsto((J+1)^{-k/2}F|_{d\mid c,\ c/d\in\mathcal A})_d$,
because each state is copied at most $(J+1)^k$ times. In a fixed $d$
block the line and memory classes are orthogonal. Summing the strip
blocks gives a fixed factor by Cauchy--Schwarz, since the numbers of
adjacent source and target strips are bounded. This proves the stated
bound on the full operator for $I=S=\varnothing$.

\subsection{Steps with a return or a stored departure}

Suppose $I\ne\varnothing$. Fix $i_*=\max I$, its returning pair
$(\ell,\mathcal L_\ell)$, and the product $T_f$ of fresh primes.
Write $c'=\ell U$ with $U=dP_<T_f$. For the raw kernel, the target
coefficient of $\widehat{\mathbf w}$ is
$j_1=\Delta(\ell U-c)\equiv-\Delta c\ne0\pmod\ell$. Thus as $j_0$
varies modulo $\ell$, the target moves along an affine line that
misses the origin. Its intersection with $\mathcal L_\ell$ contains
at most one point, fixing $j_0\equiv j_*\pmod\ell$, independently of
$U$. With $\rho_*=\tau_\rho(\widehat{\mathbf w})/\tau_\rho(\mathbf w)$,
the first-coordinate ratio condition then gives
\[
 \Delta U\rho_*\pmod1\ \in\
 \frac{[1/16,16]+\Delta c\rho_*-j_*}{\ell}\pmod1.
\]
This arc has length less than $16e^{-\lambda_{i_*}}$. For a
nonexceptional $T_f$, the second good-point test permits at most one
$U$. The edge condition $(a,c)=1$ gives $d=(c,U)$, recovering $d$,
$b$, and $P_<$, and each return then has a unique memory slot.
The fixed return pair thus gains $(J+1)^{-k}$. There are at most $M_0$
choices of that pair. For exceptional $T_f$ use the coarse bound
$O_k(M_0^k)$ and total exceptional mass at most $e^{-\lambda_{i_*}/4}$,
valid with the smaller weights $\nu_i(\ell)$. The raw row bound is
\[
 O_k\bigl(M_0(J+1)^{-k}+M_0^k e^{-\lambda_1/4}\bigr).
\]
When $I=\varnothing$ the raw row bound is $O_k(1)$ by normalized
prime choices. In the reversed direction the storage and return
sets exchange roles. For $I\cup S\ne\varnothing$ at least one
direction has the saving above, and each direction is $O_k(1)$ for
the choices of $k$ below. Schur's inequality gives
$O_k(L^{1-\jmath k/2})+O_j(L^{-j})$ for any fixed $j$.

For the major kernel, fix a returning prime $\ell$, so $\ell\nmid d$.
The geometric lattice basis has determinant $d$ and is invertible
modulo $\ell$. The target line is a nonzero linear condition in
the bounded rectangle. Fixing one coordinate for which the other
has a nonzero coefficient shows that the number of possibilities
is $O(Z(\ell^{-1}+Z^{-1/10}))$. This cancels the $1/Z$ in the
kernel, with the additional exponentially small factor in $L^\alpha$.
All remaining counts have at most polynomial cost in $L$, by the
row bounds above, including $M_0^k$ and $B^3$. If storage occurs
with no return, use the reversed direction. The other direction
always admits the coarse bound. Therefore these major-kernel steps
gain every fixed logarithmic power in operator norm.

Combining the steps with $K_{\mathrm{min}}=K_{\mathrm{raw}}-K_{\mathrm{maj}}$
and the common projections shows that for
$A_0\ge G+K+\beta+4$ and $\jmath k/2>G+3$, uniformly in the fixed
real parameters and local classes,
\begin{equation}\label{eq:stepnorm}
 \|\mathsf E_j\|\le L^{-G}
\end{equation}
for all sufficiently large $N$.

\section{Collision ranks and spectral interpolation}\label{sec:collision}

Order the $n_*=k(J+1)+2kH$ potential birth slots by time. Temporarily
order the initial $J+1$ primes in each interval and give them the
individual weights $\nu_i$. Summing over their orders recovers
$(J+1)!$ for each interval. Give each unused slot its own dummy value
outside the primes. For a partition $\pi$ of the slots put
$r(\pi)=n_*-\#\pi$ and
$\mu(\pi)=\prod_{P\in\pi}(-1)^{|P|-1}(|P|-1)!$. Then
\begin{equation}\label{eq:partitionmobius}
 \mathbf1_{\text{all birth values distinct}}
 =\sum_\pi \mu(\pi)\prod_{P\in\pi}
       \mathbf1_{\text{values agree on }P}.
\end{equation}
For a fixed assignment only refinements of its equality classes
contribute. On each class the sum is the sum of permutation signs,
because a block of size $j$ becomes a cycle in $(j-1)!$ ways. This
is zero for sizes at least two and one for a singleton. The total
absolute rank-$l$ coefficient counts permutations with $n_*-l$
cycles. Its generating polynomial is $\prod_{j=0}^{n_*-1}(1+jt)$,
by inserting each new element as a singleton or into an existing
cycle, so that coefficient is at most $n_*^{2l}$.

Set $l_0=\lceil C_0(2H\log L)/\lambda_1\rceil$. In a rank-$l$
partition, each slot after the earliest in its block must take
that earliest value. Equalities involving dummy values or distinct
intervals contribute zero. Each step has at most one birth per
interval. Sequential exposure of the initial ordered slots and
the step births therefore gives one factor $a_*$ for each of the
$l$ prescribed values. Apply the absolute row bound of
Section~\ref{sec:onestep} at the $2H$ steps and the initial total
mass at most one. Each rank-$l$ term is at most $L^{2Hc_*}a_*^l$.
Since $n_*^2a_*\le e^{-\lambda_1/2}$ for large $N$, the ranks
$l\ge l_0$ contribute at most
\begin{equation}\label{eq:highrank}
 2L^{2Hc_*} e^{-\lambda_1 l_0/2}
 \le 2L^{-2H(C_0/2-c_*)}.
\end{equation}

For $l<l_0$, connect each nonsingleton block by a tree and impose
equality on each edge by circle integration of $e(\theta(v-v'))$.
Each marked slot requires a birth indicator and carries a phase
of modulus one. Average the initial marked phases over the orders
of the active products, forming a multiplier of norm at most one.
The final vector remains $b_0$. There are at most $2l$ marked
slots and affected steps. Use the coarse bound $L^{c_*}$ on those
steps and \eqref{eq:stepnorm} on the others. Taking absolute values
after these operator estimates gives the total bound
\begin{equation}\label{eq:lowrank}
 L^{-2HG}\sum_{l<l_0}n_*^{2l}L^{2l(G+c_*)}=L^{-2HG}\exp(o(H)),
\end{equation}
since $l_0\log L=O(H(\log L)^2/L^\alpha+\log L)=o(H)$.

Choose $G>E_0+3$ and $C_0/2>c_*+E_0+3$. With $A>E_0+2$ in
\eqref{eq:momentmean}, the bounds \eqref{eq:highrank} and
\eqref{eq:lowrank} give
\begin{equation}\label{eq:physicalmoment}
 0\le\mathcal T_H\ll DYX L^{-2HE_0}.
\end{equation}
The density $c(Q)$ costs at most one. The signed operator product
was retained up to the partition contribution estimates.

The coefficient bound gives $\|f_D\|_\infty\ll L^{2C}$, and
$\|f_D\|^2\ll DYX L^{4C}$. Indeed for fixed $c$ the box count is
$O(DYX/c)$, and $\kappa_J\sum_c c^{-1}\le1$ by ordering the
$J+1$ slots per interval and allowing repeats in their product sum.
For each physical root a physical endpoint of a $2H$-step path
has $Q\mid\det(\mathbf v,\mathbf v')$ and
$|\det(\mathbf v,\mathbf v')|\ll HQDZ$. There are $O(HDZ)$ possible
integer values of the determinant divided by $Q$. At fixed value
all solutions differ by integer multiples of the physically
primitive root. Its first coordinate is comparable to the
box side, giving $O(1)$ solutions per layer. Thus the symmetric
reachability relation has degree $O(HDZ)$.

Expand $f_D=\sum_{\mathbf v}f_D^\circ(\mathbf v)u_{\mathbf v}$ over
physical roots and put $g_{\mathbf v}=\|A_D^Hu_{\mathbf v}\|$.
Only reachable pairs contribute to $\|A_D^Hf_D\|^2$.
Cauchy--Schwarz and $g_{\mathbf v}g_{\mathbf v'}\le
(g_{\mathbf v}^2+g_{\mathbf v'}^2)/2$ give
$\|A_D^Hf_D\|^2\ll L^{4C}HDZ\,\mathcal T_H$.
H\"older's inequality on the spectral measure of $A_D$ at $f_D$
gives
\[
 |\langle f_D,A_Df_D\rangle|
 \le\|f_D\|^{2-1/H}\|A_D^Hf_D\|^{1/H}.
\]
Using \eqref{eq:physicalmoment} yields
\[
 D^{-1}|\langle f_D,A_D f_D\rangle|
 \ll XYL^{4C-E_0}(HDZ)^{1/(2H)}.
\]
The last factor is bounded because $\log(HDZ)\ll_k
\log L+L^{\beta+\jmath}$ and $\beta+\jmath<\varrho$. Summing the
at most $L$ values of $D$ in \eqref{eq:goodoperatorpairing}
gives, with $E_0>E+4C+3$,
\begin{equation}\label{eq:minorbound}
 |\mathcal N_Y|\ll_E XYL^{-E}.
\end{equation}

\section{Choice of parameters and completion of the proofs}\label{sec:parameters}

Take $S\ge\max(A,2K+4)$ and use the following fixed choices in the order displayed. The precision $E$ is the target for both boxed arc estimates.
\[
\begin{gathered}
 E=2(S+C_g+3K+10),\qquad C>S+C_g+12,\\
 \alpha=1/1000,\quad \beta=\jmath=1/100,\quad \varrho=1/4,\quad \zeta=(1+\beta)/2,\\
 E_0=E+4C+6,\qquad G=E_0+5,\qquad A_0=\lceil G+K+\beta+5\rceil,\\
 k=\left\lceil\frac{2(G+4)}{\jmath}\right\rceil,\qquad
 c_*=2k+3A_0+10,\qquad C_0=2(c_*+E_0+5).
\end{gathered}
\]
Use $J=\lfloor L^\jmath\rfloor$, $H=\lceil L^\varrho/2\rceil$, and $B=L^{A_0}$. These choices give $\varrho+\beta+\jmath<1$ for root averaging, $\beta+\jmath<\varrho$ for spectral interpolation, and $\jmath k/2>G+3$ for the linked steps. They also give $G>E_0+3$ and $C_0/2>c_*+E_0+3$ for the partition-rank bound. Choose the Mellin tail and rough M\"obius precision in the major-arc estimate after these parameters. Choose the Type I and the truncation precisions greater than $S+C_g+3K+10$ wherever those precisions are freely adjustable. The independent coefficient truncation is controlled by the displayed choice of $C$.

Proposition~\ref{prop:major-arcs} and \eqref{eq:minorbound} bound both $|\mathcal M_Y|$ and $|\mathcal N_Y|$ by $O_E(XYL^{-E})$. Consequently $\mathcal R_Y\ll_E XYL^{-E}$. Cauchy--Schwarz in the cofactor bounds the boxed sum by $O_E(XL^{-E/2})$. Even the weaker squared estimate $O_E(R^3Q^2L^{-E})$ suffices, since $X\ll RQ$. There are at most $R$ blocks modulo $r_1$, $O(L^2)$ dyadic boxes, $O(L)$ label scales, at most $L$ typical patterns for large $N$, and an $O(L)$ factor from Perron separation. Together with $R^{3/2}$ from the weaker square estimate, their factor is bounded by $L^{3K+8}$. Thus these contributions, the Type I sums, and the errors in Section~\ref{sec:types} are all $O_{A,K,C_g}(QL^{-S})$ per slice. The $O(QL^{4-C})$ coefficient error can be charged once per pattern, before boxed decomposition. Partial summation as in Section~\ref{sec:types} proves Theorem~\ref{thm:pair}. Summing over the slices and applying \eqref{eq:limsup} completes the proof of Theorem~\ref{thm:main}.

\section{Formal Verification Scope}\label{sec:formal-verification}
\raggedbottom
\begingroup
\raggedright

The accompanying Lean sources comprise three modules. The toolchain is
Lean~4.36.0, and the Lake project pins Mathlib to \texttt{v4.36.0}.
Rechecking each source with \verb|lake env lean -DwarningAsError=true|
produced no Lean warnings or errors. The repository audit
\path{lean/tests/RouteFAudit.lean} lists 41 Target~3A declarations.
They have no \texttt{sorry} dependencies and use only the standard axioms
\texttt{propext}, \texttt{Classical.choice}, and \texttt{Quot.sound}.

\begin{itemize}
\item \texttt{Target3AFoundation}, in namespace \texttt{Target3A},
  verifies the signed identity \texttt{physicalShift}, the residue
  calculation \texttt{productBlockModEq}, and the conclusions
  \texttt{plusShiftObstruction}, \texttt{minusShiftDivisibility},
  and \texttt{correctedCofactorEquiv}. Its determinant completion,
  lattice image, and primitive-vector statements are
  \nolinkurl{latticeMatrix_det}, \nolinkurl{latticeMatrix_range}, and
  \nolinkurl{latticeMatrix_primitive_in_box}. The last statement
  uses a prime index and the strict positive box bound.
\item \texttt{Target3AAssembly}, also in namespace \texttt{Target3A},
  verifies the prime-index congruence via
  \nolinkurl{r2_congruence_of_minus_cofactor}, the determinant
  divisibility via \nolinkurl{determinant_dvd_of_minus_cofactors},
  and the product-fibre spacing via \nolinkurl{product_fiber_spacing}.
  The swap, shear, and reversal of a matrix class are checked in
  \nolinkurl{swap_matrix_class_formula},
  \nolinkurl{swap_matrix_class_det_and_r2},
  \nolinkurl{swap_matrix_class_involutive}, and
  \nolinkurl{swap_matrix_class_modEq_iff}.
  The factor from ordered auxiliary sampling is verified in
  \nolinkurl{auxiliary_pairing_weight_identity}, and the
  prime-index parameter inequalities in
  \nolinkurl{minor_arc_parameter_budget}.
\item \texttt{CompositeTarget3A}, in the namespace of the same name,
  verifies the composite-index completion and image in
  \nolinkurl{exists_lattice_completion},
  \nolinkurl{root_coprime_of_unit},
  \nolinkurl{latticeMatrix_range_of_primitive}, and
  \nolinkurl{latticeMatrix_range}. Its
  \nolinkurl{latticeMatrix_primitive_in_box} uses the transported
  first coordinate's unit condition, and \nolinkurl{collinear_trans}
  uses the middle first coordinate's unit condition. The theorem
  \nolinkurl{kloosterman_rhs_coprime} checks the unit condition on
  the right-hand side in the completion count. Exact gcd division and
  the reduced cofactor are checked in \nolinkurl{gcd_strip},
  \nolinkurl{gcd_dvd_cofactor}, \nolinkurl{minus_cofactor_strip},
  \nolinkurl{minus_cofactor_strip_exists},
  \nolinkurl{minus_cofactor_determinant}, and
  \nolinkurl{minus_cofactor_mod_second_prime}. The theorem
  \nolinkurl{log_budget} verifies the explicit integer-power
  inequality allowing the combined loss \(R^9\).
\end{itemize}

The assembly module specifies the analytical interface through the
structure \texttt{AnalyticMinorArcInputs}. Its fields supply the
physical pairing expansion and error bounds, spectral H\"older bound,
history expansion and estimates for its three components, and scale and
geometry bounds. From these hypotheses the module derives
\nolinkurl{minor_arc_assembled_moment_bound},
\nolinkurl{minor_arc_normalized_pairing_bound},
\nolinkurl{minor_arc_physical_pairing_bound},
\nolinkurl{minor_arc_arbitrary_log_saving}, and
\nolinkurl{minor_arc_arbitrary_real_log_saving}. Thus the final
logarithmic estimate checked by the kernel is a conditional implication
whose analytic hypotheses are visible in its type. The physical operator,
history expansion and its estimates, root counting, character sums,
and K\'atai reduction are mathematical arguments in the text. The
three certificates constitute a partial formalization of the argument.
\par\endgroup

\begin{thebibliography}{99}

\bibitem{BlomerMan}
V.~Blomer and S.~H.~Man, \emph{Bounds for Kloosterman sums on $\mathrm{GL}(n)$}, Math. Ann. \textbf{390} (2024), 1171--1200.
\href{https://doi.org/10.1007/s00208-023-02777-6}{doi:10.1007/s00208-023-02777-6}.

\bibitem{JohnstonYang}
D.~R.~Johnston and A.~Yang, \emph{Some explicit estimates for the error term in the prime number theorem}, \href{https://arxiv.org/abs/2204.01980v2}{arXiv:2204.01980v2}, 2022.

\bibitem{KandhilLanguascoMoree}
N.~Kandhil, A.~Languasco, and P.~Moree, \emph{Pair correlation of zeros of Dirichlet $L$-functions: a possible path towards the conjectures of Chowla, Elliott-Halberstam and Montgomery}, Math. Ann. \textbf{394} (2026), Article~43.
\href{https://doi.org/10.1007/s00208-026-03383-y}{doi:10.1007/s00208-026-03383-y}.

\bibitem{Lichtman}
J.~D.~Lichtman, \emph{Improved bounds for the binary Chowla conjecture}, \href{https://arxiv.org/abs/2611.28392v1}{arXiv:2611.28392v1}, 2026.

\bibitem{LuWangWu}
G.~L\"u, Z.~Wang, and J.~Wu, \emph{On the Hardy--Littlewood--Chowla conjecture}, \href{https://arxiv.org/abs/2610.20242v1}{arXiv:2610.20242v1}, 2026.

\bibitem{Pascadi}
A.~Pascadi, \emph{Large sieve inequalities for exceptional Maass forms and the greatest prime factor of $n^2+1$}, Forum Math. Pi \textbf{14} (2026), e8.
\href{https://doi.org/10.1017/fmp.2026.10025}{doi:10.1017/fmp.2026.10025}.

\end{thebibliography}

\end{document}
